\documentclass[11pt,a4paper,UTF8,fontset=fandol]{ctexart}
\usepackage[top=24mm,bottom=24mm,left=24mm,right=24mm]{geometry}
\usepackage{amsmath,amssymb,mathtools}
\usepackage{longtable,booktabs,array,calc}
\usepackage{xcolor,setspace,fancyhdr,needspace,hyperref,bookmark,xurl}
\usepackage{etoolbox}
\setstretch{1.12}
\setlength{\parindent}{0pt}\setlength{\parskip}{5pt}
\setlength{\emergencystretch}{3em}\setlength{\headheight}{15pt}
\setcounter{secnumdepth}{-1}\setcounter{tocdepth}{1}
\providecommand{\tightlist}{\setlength{\itemsep}{0pt}\setlength{\parskip}{0pt}}
\pagestyle{fancy}\fancyhf{}
\fancyhead[L]{\small 随机过程五书逐题详解}
\fancyhead[R]{\small 第008批：鞅与数学金融}
\fancyfoot[C]{\thepage}
\ctexset{section={format=\Large\bfseries,beforeskip=1.4ex,afterskip=1ex},subsection={format=\large\bfseries,beforeskip=1.2ex,afterskip=.7ex}}
\AddToHook{cmd/section/before}{\Needspace{7\baselineskip}}
\AddToHook{cmd/subsection/before}{\Needspace{5\baselineskip}}
\hypersetup{hidelinks,pdftitle={Durrett第三版第5、6章习题详解：35题与105道变式},pdfauthor={随机过程五书逐题详解},pdflang={zh-CN}}
\urlstyle{same}
\begin{document}
\begin{titlepage}
\vspace*{18mm}
{\Huge\bfseries 随机过程五书\\逐题详解\par}
\vspace{14mm}
{\LARGE 第008批\par}
\vspace{8mm}
{\Large Durrett 2016第三版\\第5章：鞅\\第6章：数学金融\par}
\vspace{12mm}
{\large 35道原题\quad 105道自编变式\\全部小问、完整树、逐步证明\par}
\vfill
{\large 配套：可编辑正文、独立LaTeX、精确计算脚本\par}
\vspace{6mm}
{2026年9月7日\par}
\end{titlepage}
\tableofcontents\clearpage
\section{本批范围与版本边界}\label{ux672cux6279ux8303ux56f4ux4e0eux7248ux672cux8fb9ux754c}

本批接续第007批第4章，连续处理 Richard Durrett《Essentials of Stochastic
Processes》2016第三版
\textbf{5.1---5.10、6.1---6.25，共35道原题的全部小问}，每题另有A、B、C三道自编变式及解答，共105道。前四章不重复计入本批。下文只简述求解必要的题设，不翻印整章英文习题。

2016版第5.5节习题在印刷220---222页，第6.7节习题在247---250页。编号以这一版为准；依据可检索转录{[}S1{]}核对题号、小问，并查看作者2021稿{[}S2{]}和2011稿{[}S3{]}的相应公式、矩阵与表格。\textbf{这不是声称取得了2016年全部原始印页。}公开稿有删除、增加和重排，不能仅凭同题号认作同题。

\subsection{本批不隐藏的条件问题}\label{ux672cux6279ux4e0dux9690ux85cfux7684ux6761ux4ef6ux95eeux9898}

5.4的转录把乘积末项写成\(U_0\)，对数中的\(Y_n\)又写成\(X_n\)。本稿明确采用\(Y_0=1,Y_n=\prod_{i=1}^nU_i\)；趋于零的结论是几乎处处，不是所有可能数列都如此。

5.6的2016转录定义严格越界\(|S_n|>a\)，随后却称简单随机游走取等号。严格越界的均值是\((\lfloor a\rfloor+1)^2\)；若改为到达整数边界\(|S_n|\ge a\)，才是\(a^2\)。2021对应题确改用\(\ge\)。两种约定都证明，不暗改条件。

6.1---6.2的\(V_0\)是初始总资金。买入股票后银行余额是\(B_0=V_0-\Delta_0S_0\)，不能把\(V_0\)完整存银行又额外得到股票。

6.3的闭区间是非负鞅概率给出的上下界；三种状态都有正的实际可能性时，按通常套利定义，端点仍产生零成本且有正概率获利的组合，真正无套利集合是开区间。本稿区分``必要的闭区间界''与``完整无套利集合''。

6.7文字称call，明确支付却为\((50-S_2)^+\)；按支付式完成看跌解答，另给字面看涨解释。6.9以明确回看支付\(S_3-\min_{0\le m\le3}S_m\)为准，不再额外扣除末尾沿用的固定执行价4。6.20先定义平均价\(A_n\)，支付式却写\((S_n-4)^+\)；两种合约分别给出完整树和答案，不替作者指定唯一原意。

旧章1.9(c)非法矩阵、2.17---2.22的原版页面对照、2.58缺因子\(p\)的来源，以及4.36服务``时长/速率''分歧继续保留。其他四本教材本批没有新增解答。第6章历史报价和公司数据只作为教材数学输入，不是当前市场信息或投资建议。

\subsection{版本对应}\label{ux7248ux672cux5bf9ux5e94}

\begin{longtable}[]{@{}p{.17\linewidth}p{.52\linewidth}p{.24\linewidth}@{}}
\toprule\noalign{}
2016题号 & 对应作者稿位置 & 注意 \\
\midrule\noalign{}
\endhead
\bottomrule\noalign{}
\endlastfoot
5.1---5.2 & 2021同号，印刷182页 & 同主题 \\
5.3 & 2016印刷221页 & 不是2021新增Polya题 \\
5.4 & 2021的5.4，印刷183页 & 统一记号 \\
5.5 & 2011的5.13，印刷176页 & 2021无同题 \\
5.6---5.9 & 2021的5.5---5.8 & 5.6越界符号改变 \\
5.10 & 2021的5.10 & 分枝过程 \\
6.1---6.5 & 2021的6.5、6.6、6.1、6.2、6.3 & 顺序重排 \\
6.6---6.9 & 2021的6.7---6.10；2011同号 & 6.7名义与公式冲突 \\
6.10---6.20 & 2011同号，印刷204---205页 & 2021删除、重排 \\
6.21---6.25 & 2011同号，印刷205---206页 & 连续利率单独处理 \\
\end{longtable}

\section{预备知识：先把每个计算步骤说清楚}\label{ux9884ux5907ux77e5ux8bc6ux5148ux628aux6bcfux4e2aux8ba1ux7b97ux6b65ux9aa4ux8bf4ux6e05ux695a}

\subsection{一、鞅是条件平均不变，不只是总平均相同}\label{ux4e00ux9785ux662fux6761ux4ef6ux5e73ux5747ux4e0dux53d8ux4e0dux53eaux662fux603bux5e73ux5747ux76f8ux540c}

原书第5.1---5.2节用\(\mathcal F_n\)表示截至第\(n\)步可用的全部资料。说\(M_n\)是鞅，需要它由这些资料确定、\(E|M_n|<\infty\)，并且
\[E(M_{n+1}\mid\mathcal F_n)=M_n.\]
独立新增量\(X_{n+1}\)若均值为0，则\(S_n=\sum_{i=1}^nX_i\)满足这一等式。若各增量方差为\(\sigma^2\)，展开平方：
\[E(S_{n+1}^2\mid\mathcal F_n)=S_n^2+2S_nEX_{n+1}+EX_{n+1}^2=S_n^2+\sigma^2.\]
于是\(S_n^2-n\sigma^2\)也是鞅，这是第5.2节例5.3的平方鞅。乘积每步新因子独立且均值1时，也同样由条件期望得到乘积鞅。

\subsection{二、停止时间要先截断}\label{ux4e8cux505cux6b62ux65f6ux95f4ux8981ux5148ux622aux65ad}

停止时刻\(T\)要求到第\(n\)步就能判断是否已经停止，不能先看未来再决定过去是否结束。对\(T\wedge n=\min(T,n)\)有逐条路径恒等式
\[M_{T\wedge n}=M_0+\sum_{k=1}^n\mathbf1_{\{T\ge k\}}(M_k-M_{k-1}).\]
``第\(k\)步还继续''由前\(k-1\)步确定，故每个和项的期望都是0。因此
\[\boxed{E M_{T\wedge n}=EM_0.}\]
这就是本批使用的第5.3节定理5.10。它本身还没有给出\(EM_T=EM_0\)。

若\(T<\infty\)几乎必然，且\(|M_{T\wedge n}|\le K\)对所有\(n\)成立，则同样有\(|M_T|\le K\)，且
\[|E M_{T\wedge n}-EM_T|\le2K P(T>n)\to0.\]
这解释了第5.4节定理5.11中的有界条件。没有这个界时，需另找控制，不能直接把极限移入期望；5.7专门补出均方收敛证明。

后面还要用一个非负极限规则：若非负随机变量\(Z_n\)几乎处处趋于\(Z\)，则\(EZ\le\liminf_n EZ_n\)。这叫Fatou不等式。它为什么能成立？令\(W_m=\inf_{n\ge m}Z_n\)，那么\(W_m\)随\(m\)增加而上升到\(Z\)，且对每个\(n\ge m\)都有\(EW_m\le EZ_n\)。对非负递增变量，极限的期望等于期望的极限，即\(EW_m\uparrow EZ\)；再让\(m\)增加，就得到上述不等式。它只给一个方向，不是任意交换极限与期望。5.7先用它估计尾部平方，再补另一段论证得到真正的等号。

\subsection{三、二叉树复制只有两个线性方程}\label{ux4e09ux4e8cux53c9ux6811ux590dux5236ux53eaux6709ux4e24ux4e2aux7ebfux6027ux65b9ux7a0b}

原书第6.2节记银行每期增长因子\(R=1+r>0\)。股票现价\(S\)，下一期为\(uS,dS\)，无套利条件为\(0<d<R<u\)。持有\(\Delta\)股，银行放\(B\)，两个终值分别为
\[\Delta uS+RB=V_H,\qquad\Delta dS+RB=V_D.\] 相减后解出
\[\boxed{\Delta=\frac{V_H-V_D}{S(u-d)}.}\]
设\(p^*=(R-d)/(u-d),q^*=1-p^*\)，则\(p^*u+q^*d=R\)。将两式按此权重相加，得到
\[\boxed{V=\Delta S+B=\frac{p^*V_H+q^*V_D}{R},\qquad B=V-\Delta S.}\]
这对应式(6.5)---(6.7)及定理6.3---6.5。\(p^*\)是复制定价权重，不是对真实上涨概率的预测。\(\Delta<0\)表示卖空，\(B<0\)表示借款。所有套利与复制论证采用教材理想条件：可分割交易、可借贷与卖空、无费用、无分红、同一借贷利率。

从终端实际支付往前逐层套这个公式，就得到欧式合约的所有价值和持仓。自融资意味着每次只重新分配当前组合已有价值，不从外部再加钱。

\subsection{四、美式合约要比较立即支付和继续价值}\label{ux56dbux7f8eux5f0fux5408ux7ea6ux8981ux6bd4ux8f83ux7acbux5373ux652fux4ed8ux548cux7ee7ux7eedux4ef7ux503c}

原书第6.4节式(6.18)---(6.19)给出
\[C_n=\frac{p^*V_{n+1,H}+q^*V_{n+1,D}}R,\qquad
\boxed{V_n=\max\{g_n,C_n\}.}\]
\(g_n\)是今天立即行权所得，\(C_n\)是继续持有的价值。若\(g_n>C_n\)就行权，反之继续，相等时任一选择都好。由最后一期开始作有限归纳：下一层已经最优，现在只有这两个选择，取最大即当前全局最优。

路径记作H（上涨）、D（下跌），\(\varnothing\)是起点。HD与DH即使现价相同，也可能有不同历史平均、最高、最低值，所以路径依赖合约不能合并这两个节点。欧式表列\(S,V,\Delta,B\)；美式表列\(S,g,C,V\)和决策，另有全部终端支付。一旦实际行权就退出，不可在后续节点再次领取。

\subsection{五、连续模型的利率与正态积分}\label{ux4e94ux8fdeux7eedux6a21ux578bux7684ux5229ux7387ux4e0eux6b63ux6001ux79efux5206}

6.21---6.25的\(r\)是连续年复利率，银行因子为\(e^{rt}\)，不是离散一期的\(1+r\)。剩余年数记作\(\tau\)，标准正态密度和分布函数为
\[\varphi(z)=\frac{e^{-z^2/2}}{\sqrt{2\pi}},\qquad\Phi(x)=\int_{-\infty}^x\varphi(z)\,dz.\]
教材无分红模型的风险中性终值为
\[S_\tau=S_0e^{(r-\sigma^2/2)\tau+\sigma\sqrt\tau Z},\qquad Z\sim N(0,1).\]
令
\[d_1=\frac{\log(S_0/K)+(r+\sigma^2/2)\tau}{\sigma\sqrt\tau},\qquad d_2=d_1-\sigma\sqrt\tau.\]
把\(S_\tau>K\)转成正态不等式，对密度完成平方，便得到第6.6节公式
\[\boxed{C=S_0\Phi(d_1)-Ke^{-r\tau}\Phi(d_2),\qquad P=C+Ke^{-r\tau}-S_0.}\]
附带程序还会独立积分原始支付作数值核对。公式中的时间、利率和波动率必须使用配套单位。

\clearpage

\section{习题5.1：A的总数是鞅，固定概率怎样求}\label{ux4e60ux98985.1aux7684ux603bux6570ux662fux9785ux56faux5b9aux6982ux7387ux600eux6837ux6c42}

\textbf{书内依据：第1章1.66的六状态遗传模型；第5.2节条件期望、第5.4节定理5.11。原题印刷220页。}每位亲本的A数为0、1、2，六个无序状态按\((0,0),(0,1),(0,2),(1,1),(1,2),(2,2)\)排列。求总数鞅性及各初态最终全A的概率。

\subsection{第一步：计算下一代一个孩子的平均A数}\label{ux7b2cux4e00ux6b65ux8ba1ux7b97ux4e0bux4e00ux4ee3ux4e00ux4e2aux5b69ux5b50ux7684ux5e73ux5747aux6570}

当前双亲为\((i,j)\)，各自传出A的概率为\(i/2,j/2\)。一个孩子各取一个基因，故其A数\(K\)的条件均值为\((i+j)/2\)。下一代兄、妹各自都有这个均值，所以两人的A总数\(M_{n+1}\)满足
\[E(M_{n+1}\mid\mathcal F_n)=i+j=M_n.\]
且\(0\le M_n\le4\)，可积性自动成立。因此\(M_n\)确为鞅。

为核对完整模型，单个孩子的概率为\(p_0=(1-i/2)(1-j/2),p_2=ij/4,p_1=1-p_0-p_2\)。两个孩子条件独立，无序对\((k,l)\)在\(k<l\)时概率\(2p_kp_l\)，在\(k=l\)时概率\(p_k^2\)，因此
\[P=\begin{pmatrix}
1&0&0&0&0&0\\
1/4&1/2&0&1/4&0&0\\
0&0&0&1&0&0\\
1/16&1/4&1/8&1/4&1/4&1/16\\
0&0&0&1/4&1/2&1/4\\
0&0&0&0&0&1
\end{pmatrix}.\]

\subsection{第二步：先证明一定吸收，再停止鞅}\label{ux7b2cux4e8cux6b65ux5148ux8bc1ux660eux4e00ux5b9aux5438ux6536ux518dux505cux6b62ux9785}

内部状态\((0,1),(1,2)\)一步吸收概率至少\(1/4\)，\((1,1)\)至少\(1/8\)；\((0,2)\)先必到\((1,1)\)，再一步至少\(1/8\)吸收。因此从任何内部状态，未来两步内吸收概率至少\(1/8\)。设\(T\)首次到两个端点，逐段条件化得
\[P(T>2m)\le(7/8)^m\to0.\]
所以\(T\)几乎必然有限。停止值\(M_{T\wedge n}\)都在0至4之间，可以使用有界停止定理，\(EM_T=M_0=i+j\)。最终总数只能0或4，故
\[\boxed{P_{(i,j)}(\text{全A})=(i+j)/4.}\]
六个答案依次为\(\boxed{0,1/4,1/2,1/2,3/4,1}\)。存在吸收状态本身不保证必然吸收，上面的两步估计不可省略。

\subsection{5.1-A：总数相同是否意味着下一代分布相同}\label{aux603bux6570ux76f8ux540cux662fux5426ux610fux5473ux7740ux4e0bux4e00ux4ee3ux5206ux5e03ux76f8ux540c}

\textbf{问题。}比较总数都为2的\((0,2)\)与\((1,1)\)。

\textbf{解。}前者每个孩子必有一个A，下一代总数恒为2。后者每个孩子为\(\operatorname{Bin}(2,1/2)\)，两人独立，总数为\(\operatorname{Bin}(4,1/2)\)，取0或4各有\(1/16\)概率。因此只知道当前总数2不足以统一确定下一代分布。但两种条件下均值都是2，鞅性只约束条件均值，不约束全部条件分布。

\subsection{5.1-B：不同的亲本权重}\label{bux4e0dux540cux7684ux4eb2ux672cux6743ux91cd}

\textbf{问题。}保留有序亲本标签，\(L_n=ai_n+bj_n\)何时对所有状态都是鞅？

\textbf{解。}下一代每位条件均值都是\((i+j)/2\)，所以\(E(L_{n+1}\mid i,j)=(a+b)(i+j)/2\)。要恒等于\(ai+bj\)，取\((1,0)\)得\(a=b\)；反过来相等时就是总数鞅的常数倍。因此充要条件\(\boxed{a=b}\)。

\subsection{5.1-C：最终总数的方差}\label{cux6700ux7ec8ux603bux6570ux7684ux65b9ux5dee}

\textbf{问题。}初始总数为\(m\)，求最终A总数方差。

\textbf{解。}最终取4的概率\(m/4\)，取0的概率\(1-m/4\)。所以二阶矩为\(16m/4=4m\)，均值为\(m\)，故\(\boxed{\operatorname{Var}M_T=m(4-m)}\)。初始总数2时方差4；均值保持不变并不意味着每条家系没有波动。

\bigskip

\section{习题5.2：对数正态乘积的鞅条件}\label{ux4e60ux98985.2ux5bf9ux6570ux6b63ux6001ux4e58ux79efux7684ux9785ux6761ux4ef6}

\textbf{书内依据：第5.2节例5.4乘积鞅与正态积分。原题印刷220页。}\(\eta_i\)独立同分布\(N(\mu,\sigma^2)\)，\(X_i=e^{\eta_i}\)，\(M_n=M_0\prod_{i=1}^nX_i\)，先取常数\(M_0\ne0\)。

\subsection{从密度求出一个因子的期望}\label{ux4eceux5bc6ux5ea6ux6c42ux51faux4e00ux4e2aux56e0ux5b50ux7684ux671fux671b}

当\(\sigma>0\)，
\[Ee^\eta=\frac1{\sqrt{2\pi}\sigma}\int_{-\infty}^{\infty}
\exp\left(x-\frac{(x-\mu)^2}{2\sigma^2}\right)dx.\] 将指数完成平方：
\[x-\frac{(x-\mu)^2}{2\sigma^2}
=-\frac{(x-\mu-\sigma^2)^2}{2\sigma^2}+\mu+\sigma^2/2.\]
平方部分是另一个正态密度，其积分为1，故\(Ee^\eta=e^{\mu+\sigma^2/2}\)。独立性还给\(E|M_n|=|M_0|e^{n(\mu+\sigma^2/2)}<\infty\)。

新因子独立于过去，于是
\[E(M_{n+1}\mid\mathcal F_n)=M_ne^{\mu+\sigma^2/2}.\]
因为\(M_n\)不会等于0，上式等于\(M_n\)当且仅当
\[\boxed{\mu=-\sigma^2/2.}\]
\(\sigma=0\)时因子确定为\(e^\mu\)，条件仍为\(\mu=0\)；若另允许\(M_0=0\)，过程恒为0，任意合法参数都可，不应和非零初值混淆。

\subsection{5.2-A：折现后才要求公平}\label{aux6298ux73b0ux540eux624dux8981ux6c42ux516cux5e73}

\textbf{问题。}银行每期因子为\(R>0\)，何时\(M_n/R^n\)是鞅？

\textbf{解。}每步折现乘数为\(e^{\eta_{n+1}}/R\)，其均值需等于1。因此\(\boxed{\mu+\sigma^2/2=\log R}\)。若一期长\(h\)、连续利率\(r\)，才进一步令\(R=e^{rh}\)，不能直接把\(R-1\)当成\(\log R\)。

\subsection{5.2-B：每期波动可以不同}\label{bux6bcfux671fux6ce2ux52a8ux53efux4ee5ux4e0dux540c}

\textbf{问题。}\(\eta_i\)独立但分布为\(N(\mu_i,\sigma_i^2)\)，求乘积鞅的条件。

\textbf{解。}逐步条件平均乘数为\(e^{\mu_i+\sigma_i^2/2}\)，故必须且只需每一期\(\boxed{\mu_i=-\sigma_i^2/2}\)。只要求若干期参数加起来抵消，不足以保证每个中间时刻的条件平均不变。

\subsection{5.2-C：均值不变却几乎必然趋零}\label{cux5747ux503cux4e0dux53d8ux5374ux51e0ux4e4eux5fc5ux7136ux8d8bux96f6}

\textbf{问题。}\(M_0>0,\sigma>0,\mu=-\sigma^2/2\)，求极限和方差。

\textbf{解。}大数定律给\(n^{-1}\log(M_n/M_0)=n^{-1}\sum\eta_i\to-\sigma^2/2<0\)，所以\(M_n\to0\)几乎处处。但\(EM_n=M_0\)。又\(Ee^{2\eta}=e^{2\mu+2\sigma^2}=e^{\sigma^2}\)，所以
\[\boxed{\operatorname{Var}M_n=M_0^2(e^{n\sigma^2}-1).}\]
少数巨大结果仍能承载全部期望；不能未经条件交换极限和期望。

\bigskip

\section{习题5.3：中性遗传模型的两个鞅与吸收界}\label{ux4e60ux98985.3ux4e2dux6027ux9057ux4f20ux6a21ux578bux7684ux4e24ux4e2aux9785ux4e0eux5438ux6536ux754c}

\textbf{书内依据：第1.1节例1.9，条件二项分布；第5.2节与第5.4节定理5.11。原题印刷221页。}无突变，\(X_{n+1}\mid X_n=x\sim\operatorname{Bin}(N,x/N)\)，0和\(N\)吸收。

\subsection{（a）基因数鞅与全A概率}\label{aux57faux56e0ux6570ux9785ux4e0eux5168aux6982ux7387}

二项均值为\(x\)，故\(E(X_{n+1}\mid\mathcal F_n)=X_n\)，且\(0\le X_n\le N\)。内部每个状态下一步吸收概率都正，有限多个内部状态使其有统一正下界，故吸收时间\(T=V_0\wedge V_N\)几乎必然有限。停止有界鞅得\(x=EX_T=N P_x(V_N<V_0)\)，因此
\[\boxed{P_x(V_N<V_0)=x/N.}\] 两个初始边界也分别给0和1。

\subsection{（b）多样性需要乘什么补偿因子}\label{bux591aux6837ux6027ux9700ux8981ux4e58ux4ec0ux4e48ux8865ux507fux56e0ux5b50}

令\(H(x)=x(N-x)\)。条件二阶矩为 \[E(X_{n+1}^2\mid X_n=x)=x^2+x(1-x/N).\]
于是 \[E[H(X_{n+1})\mid X_n=x]
=Nx-x^2-x+x^2/N=(1-1/N)H(x).\] 取\(N\ge2\)，记\(r=1-1/N\)，则
\[\boxed{Y_n=H(X_n)/r^n}\]
为鞅：新的条件均值恰好约掉多乘的\(r\)。每个固定\(n\)都有有限上界，故可积。

\subsection{（c）夹住未被吸收的概率}\label{cux5939ux4f4fux672aux88abux5438ux6536ux7684ux6982ux7387}

对内部整数\(x\)，\(N-1\le x(N-x)\le N^2/4\)，端点则为0。逐点有
\[(N-1)\mathbf1_{\{0<X_n<N\}}\le H(X_n)\le(N^2/4)\mathbf1_{\{0<X_n<N\}}.\]
又由（b）\(EH(X_n)=x(N-x)r^n\)。取期望、除以正数，得到
\[\boxed{\frac{4x(N-x)}{N^2}r^n\le P_x(0<X_n<N)
\le\frac{x(N-x)}{N-1}r^n.}\]
右端还可与1取最小值。奇数\(N\)可用\(\lfloor N^2/4\rfloor\)加强最大值界。\(N=1\)只有吸收状态，（b）的分母成为\(0^n\)，应另行排除，不硬套。

\subsection{5.3-A：求每一代的方差}\label{aux6c42ux6bcfux4e00ux4ee3ux7684ux65b9ux5dee}

\textbf{问题。}初值\(x\)固定，求\(\operatorname{Var}X_n\)。

\textbf{解。}\(EX_n=x\)，而\(NEX_n-EX_n^2=x(N-x)r^n\)。所以\(EX_n^2=Nx-x(N-x)r^n\)，减\(x^2\)得\(\boxed{x(N-x)(1-r^n)}\)。群体内部多样性减少，跨家系基因数的方差却增加。

\subsection{5.3-B：只有两个基因时的完整吸收分布}\label{bux53eaux6709ux4e24ux4e2aux57faux56e0ux65f6ux7684ux5b8cux6574ux5438ux6536ux5206ux5e03}

\textbf{问题。}\(N=2,X_0=1\)，求吸收时间\(T\)。

\textbf{解。}从1出发下一代为0、1、2的概率\(1/4,1/2,1/4\)。不吸收就回到同一个状态1，所以每一步重新面对吸收概率\(1/2\)。因此
\[\boxed{P(T=k)=2^{-k},\ k\ge1;\quad P(T>n)=2^{-n};\quad ET=2.}\]
一般有多个内部状态时，不能不加证明地认定吸收时间也几何。

\subsection{5.3-C：突变后如何平移缩放成鞅}\label{cux7a81ux53d8ux540eux5982ux4f55ux5e73ux79fbux7f29ux653eux6210ux9785}

\textbf{问题。}A向另一型突变概率\(u\)，反向\(v\)，\(u,v>0,u+v\ne1\)。构造一个线性变换的鞅。

\textbf{解。}下一代单次A概率为\(v+(1-u-v)X_n/N\)，条件均值\(Nv+aX_n\)，其中\(a=1-u-v\)。取\(c=Nv/(u+v)\)，有\(Nv+ac=c\)，故\(E(X_{n+1}-c\mid\mathcal F_n)=a(X_n-c)\)。因此\(\boxed{(X_n-c)/a^n}\)为鞅。\(a<0\)会交替变号，并无矛盾；\(a=0\)不可除，此时直接使用下一代条件均值恒为\(c\)。

\bigskip

\section{习题5.4：公平乘积为何几乎处处归零}\label{ux4e60ux98985.4ux516cux5e73ux4e58ux79efux4e3aux4f55ux51e0ux4e4eux5904ux5904ux5f52ux96f6}

\textbf{书内依据：第5.2节乘积鞅、第3.1节大数定律与初等积分。原题印刷221页。}\(U_i\)独立均匀\((0,1)\)，统一记\(Y_0=1,Y_n=\prod_{i=1}^nU_i\)，\(M_n=2^nY_n\)。

\subsection{（a）验证鞅}\label{aux9a8cux8bc1ux9785}

\(EU_{n+1}=1/2\)且新因子独立，所以\(E(M_{n+1}\mid\mathcal F_n)=2M_nEU_{n+1}=M_n\)。每个固定\(n\)有\(0<M_n<2^n\)，可积；故\(EM_n=1\)。

\subsection{（b）对数的长期均值}\label{bux5bf9ux6570ux7684ux957fux671fux5747ux503c}

\(\log Y_n=\sum_{i=1}^n\log U_i\)。分部积分给
\[E\log U=\int_0^1\log u\,du=[u\log u-u]_0^1=-1,\]
其中\(u\log u\to0\)；同时\(E|\log U|=1<\infty\)。大数定律适用，故\(\boxed{n^{-1}\log Y_n\to-1}\)几乎处处。原题转录此处的\(X_n\)按前述定义统一为\(Y_n\)。

\subsection{（c）从对数回到原变量}\label{cux4eceux5bf9ux6570ux56deux5230ux539fux53d8ux91cf}

\[n^{-1}\log M_n=\log2+n^{-1}\log Y_n\to\log2-1<0.\]
因此在这个极限成立的路径上，从某一步开始\(\log M_n\)被负常数乘\(n\)控制，得到
\[\boxed{M_n\longrightarrow0\quad\text{几乎处处}.}\]
它与所有\(n\)的\(EM_n=1\)同时成立；这说明交换极限与期望确实需要额外条件。

\subsection{\texorpdfstring{5.4-A：补偿系数改成\(c\)}{5.4-A：补偿系数改成c}}\label{aux8865ux507fux7cfbux6570ux6539ux6210c}

\textbf{问题。}\(Z_n=c^nY_n,c>0\)，什么时候趋于0或无穷？

\textbf{解。}\(n^{-1}\log Z_n\to\log c-1\)。所以\(c<e\)趋0、\(c>e\)趋无穷，均为几乎处处结论。\(c=e\)只知对数除\(n\)趋0，不能由此推断\(Z_n\)收敛。\(EZ_n=(c/2)^n\)，仅\(c=2\)为鞅；\(2<c<e\)时均值增长却典型路径趋零。

\subsection{5.4-B：低阶矩和高阶矩方向不同}\label{bux4f4eux9636ux77e9ux548cux9ad8ux9636ux77e9ux65b9ux5411ux4e0dux540c}

\textbf{问题。}求\(r>0\)时的\(EM_n^r\)及方差。

\textbf{解。}\(EU^r=1/(r+1)\)，独立性给
\[\boxed{EM_n^r=[2^r/(r+1)]^n,\qquad\operatorname{Var}M_n=(4/3)^n-1.}\]
\(0<r<1\)时\(x^r\)严格凹，\(E(2U)^r<(E2U)^r=1\)，阶矩衰减；\(r>1\)时严格凸，阶矩增长。不是所有衡量波动的量都必须同向变化。

\subsection{\texorpdfstring{5.4-C：曾达到\(K\)的概率上界}{5.4-C：曾达到K的概率上界}}\label{cux66feux8fbeux5230kux7684ux6982ux7387ux4e0aux754c}

\textbf{问题。}\(K>1,T_K=\inf\{n:M_n\ge K\}\)，证明曾经达到的概率不超过\(1/K\)。

\textbf{解。}在有限时刻截断，\(EM_{T_K\wedge n}=1\)。它非负，在\(T_K\le n\)上至少为\(K\)，故\(1\ge K P(T_K\le n)\)。让\(n\)增大，事件递增到\(T_K<\infty\)，得到\(\boxed{P(T_K<\infty)\le1/K}\)。无需\(ET_K\)有限，也没有将越过时的值强行当成恰好\(K\)。

\bigskip

\section{习题5.5：有利游戏的破产概率}\label{ux4e60ux98985.5ux6709ux5229ux6e38ux620fux7684ux7834ux4ea7ux6982ux7387}

\textbf{书内依据：第5.2节指数鞅、第5.4节停止及向下无跨越模型。原题印刷221页。}每次独立增加\(-1,1,2\)，各概率\(1/3\)，从整数\(i\ge1\)出发，求曾到0的概率。

\subsection{找一个公平的指数函数}\label{ux627eux4e00ux4e2aux516cux5e73ux7684ux6307ux6570ux51fdux6570}

令\(M_n=\rho^{S_n}\)。要使新乘数均值1，要求\((\rho^{-1}+\rho+\rho^2)/3=1\)，即
\[\rho^3+\rho^2-3\rho+1=(\rho-1)(\rho^2+2\rho-1)=0.\]
非平凡的\((0,1)\)内根是\(\rho=\sqrt2-1\)。

令\(T=V_0\)。向下只走1，所以破产前财富正，不能跳过0。由此\(0<\rho^{S_{T\wedge n}}\le1\)，且其期望为\(\rho^i\)。破产路径上停止值最终1；不破产路径上，原独立增量的大数定律给\(S_n/n\to2/3\)，所以\(S_n\to\infty\)，指数值趋0。统一有界可取期望极限，故
\[\boxed{P_i(T<\infty)=(\sqrt2-1)^i.}\]
\(i=0\)时已破产，概率1。正的平均收益不保证破产概率0。

\subsection{5.5-A：三种跳变概率不同}\label{aux4e09ux79cdux8df3ux53d8ux6982ux7387ux4e0dux540c}

\textbf{问题。}\(-1,1,2\)的概率\(a,b,c>0\)，和1，且\(-a+b+2c>0\)，求破产概率。

\textbf{解。}指数方程因式分解为 \[(\rho-1)[c\rho^2+(b+c)\rho-a]=0.\]
后一多项式在0负、1正、正半轴严格递增，故唯一小于1正根
\[\rho=\frac{-(b+c)+\sqrt{(b+c)^2+4ac}}{2c}.\]
重复原题有界停止与正漂移大数定律，答案\(\boxed{\rho^i}\)。正漂移条件用于保证这个根确在\((0,1)\)。

\subsection{5.5-B：把破产概率压到1\%以内}\label{bux628aux7834ux4ea7ux6982ux7387ux538bux52301ux4ee5ux5185}

\textbf{问题。}原题至少需要多少整数初始资本？

\textbf{解。}\(\rho^i\le0.01\)等价于\(i\ge\log(0.01)/\log\rho\)，因为分母为负，除法要反转不等号。该比值在5与6之间，故最小为\(\boxed6\)；可直接代入\(\rho^5>0.01,\rho^6<0.01\)验证。

\subsection{5.5-C：上方越界不能忽略超越量}\label{cux4e0aux65b9ux8d8aux754cux4e0dux80fdux5ffdux7565ux8d85ux8d8aux91cf}

\textbf{问题。}从1出发，第一次到0或至少2就结束，求先破产概率并检查指数等式。

\textbf{解。}第一步已经结束，0、2、3各概率\(1/3\)。所以破产概率\(\boxed{1/3}\)，且\(E\rho^{S_T}=(1+\rho^2+\rho^3)/3=\rho\)。若把2、3都当作终点2，就会得到错误方程\(\rho=h+(1-h)\rho^2\)；向上可跳2，必须保留两种上出口。

\bigskip

\section{习题5.6：退出时间下界与严格越界的等号问题}\label{ux4e60ux98985.6ux9000ux51faux65f6ux95f4ux4e0bux754cux4e0eux4e25ux683cux8d8aux754cux7684ux7b49ux53f7ux95eeux9898}

\textbf{书内依据：第5.2节平方鞅，第5.3节定理5.10。原题印刷221页。}独立增量均值0、共同方差\(\sigma^2>0\)，\(S_0=0\)。2016转录取\(T=\inf\{n:|S_n|>a\}\)，\(a\ge0\)。

\subsection{先在有限时刻证明下界}\label{ux5148ux5728ux6709ux9650ux65f6ux523bux8bc1ux660eux4e0bux754c}

平方鞅给\(ES_{T\wedge n}^2=\sigma^2E(T\wedge n)\)。如果\(ET=\infty\)，下界已经成立；否则\(T\)几乎必然有限。在\(T\le n\)上，停止位置绝对值超过\(a\)，故
\[\sigma^2E(T\wedge n)\ge a^2P(T\le n).\]
左边递增趋于\(\sigma^2ET\)，右边趋于\(a^2\)，得到\(\boxed{ET\ge a^2/\sigma^2}\)。这一步不需要停止后的平方统一有界。

\subsection{简单游走在原严格符号下并不取题称的等号}\label{ux7b80ux5355ux6e38ux8d70ux5728ux539fux4e25ux683cux7b26ux53f7ux4e0bux5e76ux4e0dux53d6ux9898ux79f0ux7684ux7b49ux53f7}

对公平\(\pm1\)游走，严格越过\(a\)实际首先碰到整数边界\(b=\lfloor a\rfloor+1\)。所有\(S_{T\wedge n}\)都在\([-b,b]\)，因此\(E(T\wedge n)\le b^2\)，先取极限得到\(ET\le b^2<\infty\)。再用有界停止位置平方趋于\(b^2\)，得
\[\boxed{ET=(\lfloor a\rfloor+1)^2\quad\text{（严格越界）}.}\]
\(a=1\)时平均4步，不是1步。因此2016转录的严格符号与后文等号声称不相容。

若明确改为\(|S_n|\ge a\)且\(a\)为正整数，停止位置\(\pm a\)，同一证明给\(ET=a^2\)；非整数\(a\)则为\(\lceil a\rceil^2\)。2021对应题采用\(\ge\)。不能把改动后的等号当作严格越界原式的证明。

\subsection{5.6-A：有界增量也能给出上界}\label{aux6709ux754cux589eux91cfux4e5fux80fdux7ed9ux51faux4e0aux754c}

\textbf{问题。}再有\(|X_i|\le c\)，求\(ET\)上界。

\textbf{解。}停止前绝对值至多\(a\)，最后一步至多\(c\)，所以\(|S_{T\wedge n}|\le a+c\)。平方等式给\(\sigma^2E(T\wedge n)\le(a+c)^2\)，取极限得
\[\boxed{a^2/\sigma^2\le ET\le(a+c)^2/\sigma^2.}\]
有限期等式先保证均值有限，不循环地提前使用无限停止结论。

\subsection{5.6-B：左右边界不对称}\label{bux5de6ux53f3ux8fb9ux754cux4e0dux5bf9ux79f0}

\textbf{问题。}公平简单游走从0出发，首次到\(-a\)或\(b\)，其中\(a,b\)正整数。求右出口概率和平均时间。

\textbf{解。}在有限区间每隔\(a+b\)步都有统一正概率离开，故停止几乎必然有限。停止线性鞅，\(0=-a(1-h)+bh\)，得\(h=a/(a+b)\)。停止平方鞅得
\[ET=a^2\frac b{a+b}+b^2\frac a{a+b}=\boxed{ab}.\]
取\(a=b\)才回到平方公式。

\subsection{5.6-C：原地停留会延长多少}\label{cux539fux5730ux505cux7559ux4f1aux5ef6ux957fux591aux5c11}

\textbf{问题。}每步以\(r/2\)左右移动、以\(1-r\)停留，\(0<r\le1\)，严格越过\(a\)的均值？

\textbf{解。}单步方差\(r\)，实际停止边界仍\(b=\lfloor a\rfloor+1\)。同样的停止平方证明得\(\boxed{ET=b^2/r}\)。也可先数平均\(b^2\)次实际移动，每次移动等均值\(1/r\)的几何日历步。\(r=0\)时始终不动，\(T=\infty\)，不可代入含\(1/r\)的公式。

\bigskip

\section{习题5.7：Wald第二等式的完整极限证明}\label{ux4e60ux98985.7waldux7b2cux4e8cux7b49ux5f0fux7684ux5b8cux6574ux6781ux9650ux8bc1ux660e}

\textbf{书内依据：第5.2节平方鞅、第5.3节有限停止；附录非负极限与Cauchy--Schwarz。原题印刷221页。}独立\(X_i\)均值0、方差\(\sigma^2<\infty\)，\(S_0=0\)，\(T\)为自然信息下的停止时刻且\(ET<\infty\)。证明\(ES_T^2=\sigma^2ET\)。

\subsection{第一步：把停止指标放入随机和}\label{ux7b2cux4e00ux6b65ux628aux505cux6b62ux6307ux6807ux653eux5165ux968fux673aux548c}

令\(D_i=X_i\mathbf1_{\{T\ge i\}}\)。``第\(i\)步还继续''在前\(i-1\)步已确定，与\(X_i\)独立，所以\(ED_i^2=\sigma^2P(T\ge i)\)。对\(i<j\)，\(D_i\mathbf1_{\{T\ge j\}}\)在第\(j-1\)步已知，而\(E(X_j\mid\mathcal F_{j-1})=0\)，因此\(E(D_iD_j)=0\)。

这证明的是正交，即交叉乘积期望0；不必声称这些带停止指标的变量独立。

\subsection{第二步：估计截断后的尾部平方}\label{ux7b2cux4e8cux6b65ux4f30ux8ba1ux622aux65adux540eux7684ux5c3eux90e8ux5e73ux65b9}

因为\(S_{T\wedge n}=\sum_{i=1}^nD_i\)，展开平方得
\[ES_{T\wedge n}^2=\sigma^2\sum_{i=1}^nP(T\ge i)=\sigma^2E(T\wedge n).\]
对\(m>n\)同样有
\[E(S_{T\wedge m}-S_{T\wedge n})^2=\sigma^2\sum_{i=n+1}^mP(T\ge i).\]
\(ET<\infty\)使\(T\)几乎必然有限，故\(m\to\infty\)时括号内趋于\(S_T-S_{T\wedge n}\)。对非负平方使用Fatou不等式得到
\[E(S_T-S_{T\wedge n})^2
\le\sigma^2\sum_{i=n+1}^{\infty}P(T\ge i)
=\sigma^2E(T-n)^+\longrightarrow0.\]
最后是收敛尾和趋零。取\(n=0\)还知\(ES_T^2\le\sigma^2ET<\infty\)。

\subsection{第三步：均方接近使二阶矩接近}\label{ux7b2cux4e09ux6b65ux5747ux65b9ux63a5ux8fd1ux4f7fux4e8cux9636ux77e9ux63a5ux8fd1}

记\(A_n=S_{T\wedge n},A=S_T\)。平方差分解后
\[|EA_n^2-EA^2|\le\sqrt{E(A_n-A)^2}\sqrt{E(A_n+A)^2}.\]
第一因子趋零，第二因子平方至多\(2EA_n^2+2EA^2\le4\sigma^2ET\)。因此\(EA_n^2\to EA^2\)。有限停止等式两边取极限，终于得到
\[\boxed{ES_T^2=\sigma^2ET.}\]
直接写\(E(S_T^2-\sigma^2T)=0\)会把这段需要证明的极限偷偷当作已知。

\subsection{5.7-A：投注大小由过去决定}\label{aux6295ux6ce8ux5927ux5c0fux7531ux8fc7ux53bbux51b3ux5b9a}

\textbf{问题。}\(H_i\)由前\(i-1\)步资料确定，\(|H_i|\le C\)，求\(Z=\sum_{i=1}^TH_iX_i\)的二阶矩。

\textbf{解。}把上面的\(D_i\)替换为\(H_iX_i\mathbf1_{\{T\ge i\}}\)，交叉项仍为0，平方项为\(\sigma^2E[H_i^2\mathbf1_{\{T\ge i\}}]\)。其尾和受\(C^2ET\)控制，同一均方收敛证明给
\[\boxed{EZ^2=\sigma^2E\sum_{i=1}^TH_i^2.}\]
投注可随过去变化，但不能先知道本步输赢。

\subsection{5.7-B：偷看未来的反例}\label{bux5077ux770bux672aux6765ux7684ux53cdux4f8b}

\textbf{问题。}独立公平\(X_1,X_2\in\{-1,1\}\)，相同就取\(T=1\)，不同取\(T=2\)。检查公式。

\textbf{解。}\(ET=(1+2)/2=3/2\)；相同时停止平方1，不同时停止和0，所以\(ES_T^2=1/2\)。二者不等。这个\(T\)不是停止时刻，因为第1步时无法判断它是否为1；方差公式的停止条件不可删除。

\subsection{5.7-C：几乎必然停止仍不足够}\label{cux51e0ux4e4eux5fc5ux7136ux505cux6b62ux4ecdux4e0dux8db3ux591f}

\textbf{问题。}公平游走从0首次到1停止，为什么不能推断\(ET=1\)？

\textbf{解。}先在\(-m\)或1停止，先到1概率\(m/(m+1)\)。令\(m\to\infty\)，知最终到1概率1，所以\(S_T=1\)几乎必然。若\(ET\)有限，Wald一式会给\(ES_T=0\)，与1矛盾，故\(ET=\infty\)。此时\(ES_T^2=1\)，不等于\(\sigma^2ET=\infty\)，正说明有限均值条件不可省。

\bigskip

\section{习题5.8：破产时间的方差与资本的线性关系}\label{ux4e60ux98985.8ux7834ux4ea7ux65f6ux95f4ux7684ux65b9ux5deeux4e0eux8d44ux672cux7684ux7ebfux6027ux5173ux7cfb}

\textbf{书内依据：第5.4.2节Wald一式、式(5.13)和习题5.7。原题印刷221页。}向上1概率\(p<1/2\)，向下1概率\(q=1-p\)；从整数\(x\ge0\)开始，\(T=V_0\)，记\(d=q-p>0\)。

\subsection{（a）先证明所需矩存在，再求方差}\label{aux5148ux8bc1ux660eux6240ux9700ux77e9ux5b58ux5728ux518dux6c42ux65b9ux5dee}

停止线性鞅给\(E_xS_{T\wedge n}=x-dE_x(T\wedge n)\ge0\)，所以\(E_xT\le x/d<\infty\)。Wald一式于是可合法应用：\(0-x=-dE_xT\)，得\(ET=x/d\)。

将增量中心化\(Y_i=X_i+d\)，其均值0、方差\(4pq\)。到停止时
\[\sum_{i=1}^TY_i=S_T-x+dT=dT-x.\] 由5.7， \[E(dT-x)^2=4pq ET=4pq x/d.\]
右边有限也保证\(ET^2<\infty\)，并未先假设未知方差存在。又\(x=dET\)，左边是\(d^2\operatorname{Var}T\)，故
\[\boxed{\operatorname{Var}_xT=\frac{4pq x}{(q-p)^3}
=\frac{x[1-(p-q)^2]}{(q-p)^3}.}\]
\(p=0\)时\(T=x\)确定，方差0；\(x=0\)时\(T=0\)。

\subsection{\texorpdfstring{（b）为什么必是常数乘\(x\)}{（b）为什么必是常数乘x}}\label{bux4e3aux4ec0ux4e48ux5fc5ux662fux5e38ux6570ux4e58x}

每次向下只能走1，故必须依次首次到达\(x-1,x-2,\ldots,0\)。每次到新层后，尚未用到的独立增量重新开始。可以先按``上一次到达发生于确定时刻\(n\)''条件化，未来增量独立且分布不随\(n\)改变，再对\(n\)求和，这就逐段证明各下降一层所需时间独立同分布。

因此\(T=T_1+\cdots+T_x\)，每段分布等同从1到0的时间。各段方差有限，相加便得\(\operatorname{Var}_xT=x\operatorname{Var}_1T\)。

\subsection{5.8-A：算出一个完整例子}\label{aux7b97ux51faux4e00ux4e2aux5b8cux6574ux4f8bux5b50}

\textbf{问题。}\(p=1/3,x=3\)，求均值、方差、二阶矩。

\textbf{解。}\(d=1/3\)，所以\(ET=9\)，方差\(4(1/3)(2/3)3/(1/3)^3=72\)；二阶矩再加均值平方，得\(\boxed{ET=9,\operatorname{Var}T=72,ET^2=153}\)。标准差为\(6\sqrt2\)，不是72。

\subsection{5.8-B：初始资本也随机}\label{bux521dux59cbux8d44ux672cux4e5fux968fux673a}

\textbf{问题。}独立初值\(J\)为非负整数，均值\(m\)、方差\(v\)，求\(T\)的均值方差。

\textbf{解。}条件于\(J=j\)，均值\(j/d\)、方差\(4pqj/d^3\)。全期望、全方差得
\[\boxed{ET=m/d,\qquad\operatorname{Var}T=4pqm/d^3+v/d^2.}\]
把\(J\)直接换成\(m\)会漏掉初始资本波动的第二项。

\subsection{5.8-C：资本大时相对波动有多小}\label{cux8d44ux672cux5927ux65f6ux76f8ux5bf9ux6ce2ux52a8ux6709ux591aux5c0f}

\textbf{问题。}证明\(T/x\)依概率趋于\(1/d\)并给定量界。

\textbf{解。}均值\(E(T/x)=1/d\)，方差\(4pq/(xd^3)\)。切比雪夫给
\[\boxed{P_x(|T/x-1/d|>\varepsilon)\le4pq/(xd^3\varepsilon^2)\to0.}\]
绝对时间仍随机，只是相对于\(x\)的误差缩小。

\bigskip

\section{习题5.9：破产时间的生成函数与选根}\label{ux4e60ux98985.9ux7834ux4ea7ux65f6ux95f4ux7684ux751fux6210ux51fdux6570ux4e0eux9009ux6839}

\textbf{书内依据：第5.2节例5.5指数鞅，第5.3---5.4节停止。原题印刷221---222页。}沿用\(0<p<1/2,q=1-p\)，初值\(x\)，\(T=V_0\)，定义\(\phi(\theta)=pe^\theta+qe^{-\theta}\)。

\subsection{（a）停止指数鞅}\label{aux505cux6b62ux6307ux6570ux9785}

\(M_n=e^{\theta S_n}/\phi(\theta)^n\)为鞅。对\(\theta<0\)，\(\phi'(\theta)=pe^\theta-qe^{-\theta}<0\)，而\(\phi(0)=1\)，所以\(\phi(\theta)>1\)。停止前后\(S_{T\wedge n}\ge0\)，因此\(0<M_{T\wedge n}\le1\)。由5.8，\(T\)几乎必然有限，停止时\(S_T=0\)，故有界收敛给
\[\boxed{E_x[\phi(\theta)^{-T}]=e^{\theta x},\qquad\theta\le0.}\]
\(\theta=0\)时两边都是1。

\subsection{\texorpdfstring{（b）改用参数\(s\)并选择正确根}{（b）改用参数s并选择正确根}}\label{bux6539ux7528ux53c2ux6570sux5e76ux9009ux62e9ux6b63ux786eux6839}

给定\(0<s<1\)，令\(\phi(\theta)=1/s\)。写\(z=e^\theta\in(0,1)\)，方程为\(ps z^2-z+qs=0\)。在0处为正，在1处为\(s-1<0\)，故较小根在\((0,1)\)，另一根大于1。于是
\[z=\frac{1-\sqrt{1-4pqs^2}}{2ps},\qquad
\boxed{E_xs^T=f(s)^x,\quad f(s)=\frac{1-\sqrt{1-4pqs^2}}{2ps}.}\]
\(s=1\)取极限为1，\(s=0\)时\(x>0\)取0，\(x=0\)时\(T=0\)故恒1。\(p=0\)不可除，直接由\(T=x\)得到\(s^x\)。

\subsection{（c）幂形式的独立解释}\label{cux5e42ux5f62ux5f0fux7684ux72ecux7acbux89e3ux91ca}

由5.8的逐层下降分解，\(T\)是\(x\)个独立同分布跨层时间之和。\(s\)的指数相加变成乘法，再用独立性，\(Es^{T_1+\cdots+T_x}=\prod Es^{T_i}=f(s)^x\)。\(x=0\)是空和0、空乘积1。

\subsection{5.9-A：有利游走条件在最终破产}\label{aux6709ux5229ux6e38ux8d70ux6761ux4ef6ux5728ux6700ux7ec8ux7834ux4ea7}

\textbf{问题。}改成\(p>1/2\)，条件在最终破产后，破产时间的生成函数是什么？

\textbf{解。}指数停止给破产概率\(h(i)=(q/p)^i\)。条件于将来破产，当前\(i\)下一步向上概率变为\(p h(i+1)/h(i)=q\)，向下为\(q h(i-1)/h(i)=p\)。上下概率互换成为向下漂移，故
\[\boxed{E_x(s^T\mid T<\infty)=\left[\frac{1-\sqrt{1-4pqs^2}}{2qs}\right]^x.}\]
先改变条件路径分布，再应用生成函数，不能把原有利漂移直接放进有限均值公式。

\subsection{5.9-B：最早几种破产时间}\label{bux6700ux65e9ux51e0ux79cdux7834ux4ea7ux65f6ux95f4}

\textbf{问题。}从1出发求\(P(T=1),P(T=3),P(T=5)\)，解释偶数为何不可能。

\textbf{解。}净位移为\(-1\)，与步数同奇偶，所以只能奇数。第一步向下概率\(q\)；三步首次到0唯一序列\(+--\)，概率\(pq^2\)；五步有\(++---\)和\(+-+--\)两条，各\(p^2q^3\)。故\(\boxed{q,pq^2,2p^2q^3}\)，也对应生成函数最早三项。仅计算最终在0会错误地纳入此前已破产的路径。

\subsection{5.9-C：加入原地停留的日历时间}\label{cux52a0ux5165ux539fux5730ux505cux7559ux7684ux65e5ux5386ux65f6ux95f4}

\textbf{问题。}每步以\(h\in[0,1)\)停留，否则按\(p,q\)实际移动，求日历破产时间生成函数。

\textbf{解。}一次实际移动前的日历等待\(G\)几何分布，\(P(G=k)=(1-h)h^{k-1}\)，所以\(Es^G=(1-h)s/(1-hs)\)。条件于所需实际步数\(m\)，总日历时间是\(m\)个独立\(G\)之和，再对\(m\)平均得到
\[\boxed{E_xs^{T_{\rm clock}}=\left[f\left(\frac{(1-h)s}{1-hs}\right)\right]^x.}\]
均值为\(x/[(q-p)(1-h)]\)；空间移动次数与钟表步数不同。

\bigskip

\section{习题5.10：分枝过程的指数函数与灭绝概率}\label{ux4e60ux98985.10ux5206ux679dux8fc7ux7a0bux7684ux6307ux6570ux51fdux6570ux4e0eux706dux7eddux6982ux7387}

\textbf{书内依据：第1.11节分枝过程、第5.2节条件期望与第5.4节极限。原题印刷222页。}后代数\(K\)分布\(p_j\)，\(p_0>0,\mu=EK>1\)，不同个体、世代独立。\(\phi(s)=\sum p_js^j\)，\(0\le s\le1\)，约定\(s^0=1\)。

\subsection{（a）直接计算条件生成函数}\label{aux76f4ux63a5ux8ba1ux7b97ux6761ux4ef6ux751fux6210ux51fdux6570}

给定\(Z_n=z\)，下一代为\(K_1+\cdots+K_z\)，独立性使
\[E(s^{Z_{n+1}}\mid Z_n=z)=E\prod_{i=1}^zs^{K_i}=\prod_{i=1}^zEs^{K_i}=\phi(s)^z.\]
\(z=0\)时空乘积也为1，因此
\[\boxed{E(s^{Z_{n+1}}\mid\mathcal F_n)=\phi(s)^{Z_n}.}\]

\subsection{（b）小于1的不动点为什么是灭绝概率}\label{bux5c0fux4e8e1ux7684ux4e0dux52a8ux70b9ux4e3aux4ec0ux4e48ux662fux706dux7eddux6982ux7387}

\(\phi(0)=p_0>0\)，\(\phi(1)=1\)，且\(\phi(s)-s\)在1左侧充分近时为负。有限均值时由左导数\(\mu-1>0\)得到；若允许无限均值，也可由\([1-\phi(s)]/(1-s)=E(1+\cdots+s^{K-1})\uparrow EK>1\)得到。存在某种\(j\ge2\)的正概率，因此\(\phi(s)-s\)在\((0,1)\)严格凸，从而恰有一个根\(\rho\in(0,1)\)，边界另有根1。由（a），\(\rho^{Z_n}\)是有界鞅。

从一人开始，令\(q_n=P(Z_n=0)\)。\(q_0=0,q_{n+1}=\phi(q_n)\)，灭绝事件逐代递增，所以\(q_n\uparrow q\)。连续性给\(q=\phi(q)\)；又\(q_0\le\rho\)，递推得每个\(q_n\le\rho\)，故\(q\le\rho\)。\(\rho\)为最小非负根，遂\(q=\rho\)。

鞅的极限解释也可直接补足：若生存路径无限次进入\(1,\ldots,L\)，每次下一代全灭的条件概率至少\(p_0^L>0\)，连续躲过\(m\)次这类机会的概率至多\((1-p_0^L)^m\)。因此生存时\(Z_n\to\infty\)，\(\rho^{Z_n}\to0\)；灭绝时最终为1。有界期望极限给灭绝概率\(\rho^{Z_0}\)。

从\(k\)人开始有\(k\)个独立支系，必须都灭绝，所以\(\boxed{P_k(T_0<\infty)=\rho^k}\)。\(k=0\)时本已灭绝，概率1。

\subsection{5.10-A：只有0或2个后代}\label{aux53eaux67090ux62162ux4e2aux540eux4ee3}

\textbf{问题。}\(p_0=1/4,p_2=3/4\)，从\(k\)人开始求灭绝概率。

\textbf{解。}\(\phi(s)=1/4+3s^2/4\)，不动点方程\((3s-1)(s-1)=0\)，较小根\(1/3\)。因此\(\boxed{P_k(\text{灭绝})=3^{-k}}\)。平均后代\(3/2>1\)并不排除灭绝。

\subsection{5.10-B：已知最终灭绝，首代分布怎样改变}\label{bux5df2ux77e5ux6700ux7ec8ux706dux7eddux9996ux4ee3ux5206ux5e03ux600eux6837ux6539ux53d8}

\textbf{问题。}上一模型，条件在最终灭绝后，求首位后代分布及全家系总人数期望。

\textbf{解。}若有\(j\)名后代，后续全灭概率\(\rho^j\)，所以条件后代概率\(\widetilde p_j=p_j\rho^{j-1}\)。代入得\(\widetilde p_0=3/4,\widetilde p_2=1/4\)，均值\(1/2\)。条件后独立子家系仍各自按灭绝条件独立，第\(n\)代均值\((1/2)^n\)。包含始祖的总人数期望为\(\sum_{n\ge0}2^{-n}=\boxed2\)；\(k\)条家系则\(2k\)。

\subsection{5.10-C：始祖数也随机}\label{cux59cbux7956ux6570ux4e5fux968fux673a}

\textbf{问题。}\(J\sim\operatorname{Pois}(a)\)且与生殖独立，求全体灭绝概率。

\textbf{解。}给定\(J=j\)为\(\rho^j\)，再平均：
\[\boxed{P(\text{灭绝})=e^{-a}\sum_{j=0}^{\infty}\frac{(a\rho)^j}{j!}=e^{-a(1-\rho)}.}\]
\(j=0\)的空家系也属于灭绝事件，已由和式第一项包含。

\clearpage

\section{习题6.1：一时期看涨期权的价格与复制}\label{ux4e60ux98986.1ux4e00ux65f6ux671fux770bux6da8ux671fux6743ux7684ux4ef7ux683cux4e0eux590dux5236}

\textbf{书内依据：第6.2节式(6.5)---(6.7)。原题印刷247页。}股票初价110，下一期121或99，利率4\%，执行价113。

\subsection{（a）初始价值}\label{aux521dux59cbux4ef7ux503c}

风险中性平均要求\(121p+99(1-p)=110(26/25)=114.4\)，故\(p=7/10\)。两个支付为8、0，所以
\[\boxed{V_0=\frac{(7/10)8}{26/25}=70/13.}\] \#\# （b）股票份额

两个支付差8，两个股票终价差22，故\(\boxed{\Delta_0=4/11}\)股。 \#\#
（c）现金与股票逐状态验证

买股需40，初始总资金只有\(70/13\)，银行余额因此为\(B_0=70/13-40=-450/13\)，负数表示借款。下一期银行为\(-36\)；上涨股票终值44，净值8；下跌股票终值36，净值0，恰好复制。不能完整保留\(V_0\)现金又另外买股票，否则初始成本重复计算。

\subsection{6.1-A：同执行价看跌}\label{aux540cux6267ux884cux4ef7ux770bux8dcc}

\textbf{问题。}求\((113-S_1)^+\)的价与组合。

\textbf{解。}支付为0、14，价格\(\boxed{105/26}\)，份额\(\Delta=-14/22=-7/11\)，银行\(B=105/26+70=1925/26\)。下一期银行77，股票为\(-77\)或\(-63\)，合计0、14。看跌用卖空股票复制。

\subsection{6.1-B：上涨达标只付1元}\label{bux4e0aux6da8ux8fbeux6807ux53eaux4ed81ux5143}

\textbf{问题。}求\(\mathbf1_{\{S_1>113\}}\)价格与份额。

\textbf{解。}支付1、0，价格\(\boxed{35/52}\)，份额\(1/22\)，银行\(35/52-5=-225/52\)。下一期银行\(-9/2\)，股票为\(11/2\)或\(9/2\)，相加为1、0，完成复制。

\subsection{6.1-C：一个偏高报价的模型反例}\label{cux4e00ux4e2aux504fux9ad8ux62a5ux4ef7ux7684ux6a21ux578bux53cdux4f8b}

\textbf{问题。}合约报价6，允许教材理想借贷卖空，证明不相容。

\textbf{解。}卖合约收6，买复制组合花\(70/13\)，初始多\(8/13\)，把它存银行。到期复制组合恰好抵消卖出合约的全部负债，另得\((8/13)(26/25)=\boxed{16/25}\)，每个状态都正。由零外部投入建立的这个组合证明理想模型中的套利，不是现实市场交易建议。

\bigskip

\section{习题6.2：一时期看跌的负股票份额}\label{ux4e60ux98986.2ux4e00ux65f6ux671fux770bux8dccux7684ux8d1fux80a1ux7968ux4efdux989d}

\textbf{书内依据：第6.2节二状态复制。原题印刷247页。}股票60，下一期75或45，利率5\%，支付\((60-S_1)^+\)。

\subsection{（a）价格}\label{aux4ef7ux683c}

\(75p+45(1-p)=63\)，故\(p=3/5\)。支付为0、15，所以\(\boxed{V_0=(2/5)15/(21/20)=40/7}\)。
\#\# （b）需要卖出多少股票

\(\boxed{\Delta_0=(0-15)/(75-45)=-1/2}\)，负号表示卖空半股。 \#\#
（c）检查到期财富

银行\(B_0=40/7+30=250/7\)，下一期变为\(75/2\)；股票头寸为\(-75/2\)或\(-45/2\)，相加0或15。因此\(\boxed{(V_0,\Delta_0,B_0)=(40/7,-1/2,250/7)}\)。银行资金包括卖空股票收到的款项。

\subsection{6.2-A：由看跌得到看涨}\label{aux7531ux770bux8dccux5f97ux5230ux770bux6da8}

\textbf{问题。}求相同执行价的看涨价格。

\textbf{解。}逐状态\(C_1-P_1=S_1-60\)，折现期望给\(C_0-P_0=60-60/R\)，所以\(\boxed{C_0=60/7}\)。直接用上涨概率计算\((3/5)15/(21/20)\)也同值。

\subsection{6.2-B：增加现在行权的选择}\label{bux589eux52a0ux73b0ux5728ux884cux6743ux7684ux9009ux62e9}

\textbf{问题。}0或1期可行权的美式看跌价值是多少？

\textbf{解。}立即支付\((60-60)^+=0\)，继续值\(40/7>0\)，所以继续，美式价值仍\(\boxed{40/7}\)。增加权利只保证价格不降，不保证严格增加。

\subsection{6.2-C：银行利率比两种股票增长都高}\label{cux94f6ux884cux5229ux7387ux6bd4ux4e24ux79cdux80a1ux7968ux589eux957fux90fdux9ad8}

\textbf{问题。}若银行因子改为\(13/10\)，是否还能定价？

\textbf{解。}卖空一股收到60存银行，到期78，买回股票只需75或45，净赚3或33。基础股票---银行市场本身已套利。算出的\(p=(78-45)/30=11/10>1\)不是概率，不能继续拿它给期权定价。

\bigskip

\section{习题6.3：三个状态、价格上下界与端点}\label{ux4e60ux98986.3ux4e09ux4e2aux72b6ux6001ux4ef7ux683cux4e0aux4e0bux754cux4e0eux7aefux70b9}

\textbf{书内依据：第6.1节无套利比较，第6.2节鞅概率与线性方程。原题印刷247页。}利率0，股票100，终价130、110、80；期权支付\((S_1-105)^+\)即25、5、0。

\subsection{（a）所有非负鞅概率}\label{aux6240ux6709ux975eux8d1fux9785ux6982ux7387}

设\(p_1=t\)，归一化与股票均值要求给
\[p_1+p_2+p_3=1,\quad130p_1+110p_2+80p_3=100.\]
消去\(p_3\)，有\(50t+30p_2=20\)，所以
\[\boxed{p_1=t,\ p_2=2/3-5t/3,\ p_3=1/3+2t/3,\quad0\le t\le2/5.}\]
区间两端分别来自\(p_1,p_2\)非负。

\subsection{（b）期权期望的范围}\label{bux671fux6743ux671fux671bux7684ux8303ux56f4}

\(25p_1+5p_2=10/3+50t/3\)随\(t\)递增，因此下界\(v_0=10/3\)、上界\(v_1=10\)。

\subsection{（c）上界对应的支配组合}\label{cux4e0aux754cux5bf9ux5e94ux7684ux652fux914dux7ec4ux5408}

以总资金10持股\(\boxed{x_1=1/2}\)，银行\(-40\)，终值为\((25,15,0)\)，逐状态不小于期权\((25,5,0)\)。在状态1、3取等号，状态2多10。

\subsection{（d）下界对应的被支配组合}\label{dux4e0bux754cux5bf9ux5e94ux7684ux88abux652fux914dux7ec4ux5408}

以总资金\(10/3\)持股\(\boxed{x_0=1/6}\)，银行\(-40/3\)，终值\((25/3,5,0)\)，不高于期权。在状态2、3取等号，状态1少\(50/3\)。

\subsection{（e）必要闭区间与完整无套利集合}\label{eux5fc5ux8981ux95edux533aux95f4ux4e0eux5b8cux6574ux65e0ux5957ux5229ux96c6ux5408}

若报价\(c>10\)，卖期权买上方组合，初始有正盈余，到期不亏；若\(c<10/3\)，买期权卖下方组合，同理。因此无套利价必须落在\([10/3,10]\)，这证明题面``只能位于''该区间的必要结论。

若三种状态都有正的实际可能性，端点也需排除：\(c=10\)时初始净额0，状态2仍赚10；\(c=10/3\)时状态1赚\(50/3\)，其他状态不亏。它们满足通常``零成本、绝不亏且有正概率获利''的套利定义。

任一内部价格取\(t=(3c-10)/50\in(0,2/5)\)，三个定价概率都严格正，且同时把股票、期权初价表示成期望。零成本组合的终值期望为0，不可能逐状态非负而某状态正，否则严格正概率下期望应正。故标准三状态模型中
\[\boxed{\text{真正无套利价格恰为 }(10/3,10).}\]
非负鞅概率与对每个实际可能状态严格正的概率不能混为一谈。

\subsection{6.3-A：只在状态1付1}\label{aux53eaux5728ux72b6ux60011ux4ed81}

\textbf{问题。}求该数字合约无套利价格范围。

\textbf{解。}在定价概率下价值就是\(t\)。严格正要求\(0<t<2/5\)，故\(\boxed{c\in(0,2/5)}\)。价0时免费得到有时付1的合约本身就套利，不能包含下端点。

\subsection{6.3-B：加入原期权后能确定其他合约吗}\label{bux52a0ux5165ux539fux671fux6743ux540eux80fdux786eux5b9aux5176ux4ed6ux5408ux7ea6ux5417}

\textbf{问题。}原看涨报价为内部\(c\)，求状态1数字合约的唯一价并复制。

\textbf{解。}新的价格约束唯一确定\(t=(3c-10)/50\)。逐状态恒等式
\[\mathbf1_{\{1\}}=(3/50)G-(1/100)S_1+4/5\]
在三个状态分别给1、0、0，所以它的成本\(3c/50-1+4/5=\boxed{(3c-10)/50}\)。股票、银行、这份非线性期权三列独立，足以复制任意三个终值。

\subsection{6.3-C：界限重合时为何可以取到}\label{cux754cux9650ux91cdux5408ux65f6ux4e3aux4f55ux53efux4ee5ux53d6ux5230}

\textbf{问题。}执行价改为80，求看涨价。

\textbf{解。}所有终价均不低于80，支付就等于\(S_1-80\)。一股股票加银行借款80可完全复制，初价\(\boxed{20}\)。上下界重合但并无端点套利，因为上、下组合在每个状态都和支付相同，不再存在严格的终值差额。

\bigskip

\section{习题6.4：三种赌注的套利与无套利}\label{ux4e60ux98986.4ux4e09ux79cdux8d4cux6ce8ux7684ux5957ux5229ux4e0eux65e0ux5957ux5229}

\textbf{书内依据：第6.1节无套利及线性组合。原题印刷248页。}每份初价1，三状态终值为
\[\begin{array}{c|ccc}&\text{赌注1}&\text{赌注2}&\text{赌注3}\\\hline
\text{胜}&0&2&1/2\\\text{负}&1&2&3/2\\\text{平}&3/2&0&0
\end{array}.\] 利率0，允许正负份额。

\subsection{（a）三种均可交易}\label{aux4e09ux79cdux5747ux53efux4ea4ux6613}

买1份赌注1、1份赌注2，卖1份赌注3，初始成本1，三个终值都\(3/2\)。借1元支付初始成本，到期还1，三个状态均净赚\(\boxed{1/2}\)，因此存在套利。

\subsection{（b）只有前两种}\label{bux53eaux6709ux524dux4e24ux79cd}

使两份终值均值都为1。第二份给\(2(p_{\text{胜}}+p_{\text{负}})=1\)，故\(p_{\text{平}}=1/2\)；第一份给\(p_{\text{负}}+(3/2)(1/2)=1\)，故\(p_{\text{负}}=1/4\)，剩下\(p_{\text{胜}}=1/4\)。

在严格正概率\((1/4,1/4,1/2)\)下，任意零成本的银行加两赌注组合期望终值0，不可能不亏且某状态严格获利。因此只有前两种时无套利；这是排除所有组合，不是只检查几个猜测。

\subsection{6.4-A：第三份应该多少钱}\label{aux7b2cux4e09ux4efdux5e94ux8be5ux591aux5c11ux94b1}

\textbf{问题。}保持前两份1元，求第三份相容价格。

\textbf{解。}它在上面概率下的期望是\((1/4)(1/2)+(1/4)(3/2)=\boxed{1/2}\)。也可看逐状态恒等式``赌注3=赌注1+赌注2−现金\(3/2\)''，成本亦为\(1/2\)。

\subsection{6.4-B：复制获胜数字支付}\label{bux590dux5236ux83b7ux80dcux6570ux5b57ux652fux4ed8}

\textbf{问题。}只在胜时付1，其余0，如何复制？

\textbf{解。}银行放\(3/2\)，卖1份赌注1、卖\(1/4\)份赌注2；终值分别\(3/2-0-1/2=1\)、\(3/2-1-1/2=0\)、\(3/2-3/2=0\)。初价\(3/2-1-1/4=\boxed{1/4}\)。

\subsection{6.4-C：禁止卖空赌注}\label{cux7981ux6b62ux5356ux7a7aux8d4cux6ce8}

\textbf{问题。}三种仍各价1，只能买非负份额，银行可借贷，是否仍有套利？

\textbf{解。}用（b）的严格正概率，第三份期望仅\(1/2\)，前两份各1。借款购买非负份额\(h_1,h_2,h_3\)后净利润期望为\(-h_3/2\le0\)。若终值各状态非负且某状态正，严格正概率下期望应正，矛盾。因此这种交易限制下无套利。允许卖空是（a）组合能成立的必要交易条件。

\bigskip

\section{习题6.5：两只股票复制三个终值}\label{ux4e60ux98986.5ux4e24ux53eaux80a1ux7968ux590dux5236ux4e09ux4e2aux7ec8ux503c}

\textbf{书内依据：第6.1---6.2节鞅概率、复制方程。原题印刷248页。}A初价100、三终值120、110、84；B初价50、三终值30、55、60；利率0。B看涨执行价50，支付0、5、10，要求两种求法。

\subsection{方法一：共同风险中性概率}\label{ux65b9ux6cd5ux4e00ux5171ux540cux98ceux9669ux4e2dux6027ux6982ux7387}

三个方程为
\[p_1+p_2+p_3=1,\quad120p_1+110p_2+84p_3=100,\quad30p_1+55p_2+60p_3=50.\]
消去\(p_3\)得\(36p_1+26p_2=16\)和\(30p_1+5p_2=10\)。后式\(p_2=2-6p_1\)代回得\(p_1=3/10\)，故
\[\boxed{(p_1,p_2,p_3)=(3/10,1/5,1/2),\quad V_0=5/5+10/2=6.}\]

\subsection{方法二：股票和现金组合}\label{ux65b9ux6cd5ux4e8cux80a1ux7968ux548cux73b0ux91d1ux7ec4ux5408}

银行\(b\)、A份额\(x\)、B份额\(y\)须满足
\[b+120x+30y=0,\ b+110x+55y=5,\ b+84x+60y=10.\]
后两式减第一式，\(-10x+25y=5,-36x+30y=10\)，解出
\[\boxed{x=-1/6,\quad y=2/15,\quad b=16.}\]
初始成本\(16-100/6+100/15=6\)，终值分别\(16-20+4=0\)、\(16-55/3+22/3=5\)、\(16-14+8=10\)，全部复制。

\subsection{6.5-A：只在第三状态付1}\label{aux53eaux5728ux7b2cux4e09ux72b6ux6001ux4ed81}

\textbf{问题。}求价格与组合。

\textbf{解。}价格为该定价概率\(\boxed{1/2}\)。银行\(11/2\)、A份额\(-1/24\)、B份额\(-1/60\)逐状态给0、0、1，成本\(11/2-100/24-50/60=1/2\)。

\subsection{6.5-B：B股票同执行价看跌}\label{bbux80a1ux7968ux540cux6267ux884cux4ef7ux770bux8dcc}

\textbf{问题。}求\((50-S_B)^+\)。

\textbf{解。}终值20、0、0，所以价\(20(3/10)=\boxed6\)。因为初价与执行价都50且利率0，平价也给看涨看跌同价；支付恒等式\(P=C+50-S_B\)能将原组合直接改成看跌组合。

\subsection{6.5-C：删除A股票后价格不再唯一}\label{cux5220ux9664aux80a1ux7968ux540eux4ef7ux683cux4e0dux518dux552fux4e00}

\textbf{问题。}只留B和银行，三状态仍均可能，求原看涨无套利范围。

\textbf{解。}令\(p_3=t\)，归一化与B均值给\(p_2=4/5-6t/5,p_1=1/5+t/5\)。严格正要求\(0<t<2/3\)，期权期望\(5p_2+10t=4+4t\)，故\(\boxed{c\in(4,20/3)}\)。每个内部价对应严格正概率，端点则有状态上严格的上/下支配组合，和6.3同理排除。第二只独立股票正是把原题价格唯一锁定为6的额外信息。

\bigskip

\section{习题6.6：两期看涨的每一步复制}\label{ux4e60ux98986.6ux4e24ux671fux770bux6da8ux7684ux6bcfux4e00ux6b65ux590dux5236}

\textbf{书内依据：第6.2节向后定价与自融资复制。原题印刷248页。}\(S_0=100,u=2,d=1/2,R=5/4,N=2,K=80\)，看涨支付\((S_2-80)^+\)。

\(p^*=(5/4-1/2)/(2-1/2)=1/2\)。终价400、100、100、25，支付320、20、20、0。第一期上涨节点价值\((320+20)/(2R)=136\)，下跌节点\((20+0)/(2R)=8\)。再往前得
\[\boxed{V_0=(136+8)/(2R)=288/5.}\]
起点股票份额\((136-8)/(200-50)=64/75\)，银行\(288/5-100(64/75)=-416/15\)；上涨节点持股1、银行\(-64\)，下跌节点持股\(4/15\)、银行\(-16/3\)。例如走到下跌节点，旧组合价值\((64/75)50+(5/4)(-416/15)=8\)，恰好够重新配置。以下表列出全部节点，负银行余额是借款。

\begingroup
\small
\setlength{\tabcolsep}{4pt}
\renewcommand{\arraystretch}{1.10}

\begin{longtable}[]{@{}
  >{\raggedleft\arraybackslash}p{(\columnwidth - 10\tabcolsep) * \real{0.1739}}
  >{\raggedright\arraybackslash}p{(\columnwidth - 10\tabcolsep) * \real{0.1304}}
  >{\raggedleft\arraybackslash}p{(\columnwidth - 10\tabcolsep) * \real{0.1739}}
  >{\raggedleft\arraybackslash}p{(\columnwidth - 10\tabcolsep) * \real{0.1739}}
  >{\raggedleft\arraybackslash}p{(\columnwidth - 10\tabcolsep) * \real{0.1739}}
  >{\raggedleft\arraybackslash}p{(\columnwidth - 10\tabcolsep) * \real{0.1739}}@{}}
\toprule\noalign{}
\begin{minipage}[b]{\linewidth}\raggedleft
期数
\end{minipage} & \begin{minipage}[b]{\linewidth}\raggedright
历史
\end{minipage} & \begin{minipage}[b]{\linewidth}\raggedleft
股价 \(S\)
\end{minipage} & \begin{minipage}[b]{\linewidth}\raggedleft
价值 \(V\)
\end{minipage} & \begin{minipage}[b]{\linewidth}\raggedleft
股数 \(\Delta\)
\end{minipage} & \begin{minipage}[b]{\linewidth}\raggedleft
银行 \(B\)
\end{minipage} \\
\midrule\noalign{}
\endhead
\bottomrule\noalign{}
\endlastfoot
0 & \(\varnothing\) & \(100\) & \(\frac{288}{5}\) & \(\frac{64}{75}\) &
\(-\frac{416}{15}\) \\
1 & H & \(200\) & \(136\) & \(1\) & \(-64\) \\
1 & D & \(50\) & \(8\) & \(\frac{4}{15}\) & \(-\frac{16}{3}\) \\
\end{longtable}

\textbf{全部终端历史与支付（H为涨，D为跌）。}

\begin{longtable}[]{@{}lrrr@{}}
\toprule\noalign{}
路径 & 终价 & 到期支付 & 定价概率 \\
\midrule\noalign{}
\endhead
\bottomrule\noalign{}
\endlastfoot
HH & \(400\) & \(320\) & \(\frac{1}{4}\) \\
HD & \(100\) & \(20\) & \(\frac{1}{4}\) \\
DH & \(100\) & \(20\) & \(\frac{1}{4}\) \\
DD & \(25\) & \(0\) & \(\frac{1}{4}\) \\
\end{longtable}

\endgroup

\subsection{6.6-A：同执行价看跌及其对冲}\label{aux540cux6267ux884cux4ef7ux770bux8dccux53caux5176ux5bf9ux51b2}

\textbf{问题。}求看跌价格与初始股票份额。

\textbf{解。}看跌只有DD支付55，价格\((1/4)55/R^2=\boxed{44/5}\)。由平价\(P=C+80/R^2-S_0\)同值，份额为原看涨减1，即\(\boxed{-11/75}\)；各时刻银行头寸在原看涨基础上加一份到期付80的债券。

\subsection{6.6-B：超过执行价只支付1}\label{bux8d85ux8fc7ux6267ux884cux4ef7ux53eaux652fux4ed81}

\textbf{问题。}求数字看涨初价和初始份额。

\textbf{解。}除DD外三条路径都付1，价值\((3/4)/R^2=\boxed{12/25}\)。第一期H价值\(1/R=4/5\)，D价值\(1/(2R)=2/5\)，份额差除股价差为\((2/5)/150=\boxed{1/375}\)。数字合约的对冲不是把普通看涨份额直接除以执行价。

\subsection{6.6-C：复制线性远期支付}\label{cux590dux5236ux7ebfux6027ux8fdcux671fux652fux4ed8}

\textbf{问题。}若支付\(S_2-80\)可正可负，求所有时刻的组合。

\textbf{解。}持一股，并借入到期需归还80的银行资金。第\(n\)期价值及现金为
\[\boxed{V_n=S_n-\frac{80}{R^{2-n}},\qquad\Delta_n=1,\qquad B_n=-\frac{80}{R^{2-n}}.}\]
初价\(100-80/R^2=244/5\)。线性支付允许负值，不能用看涨的正部操作替代。

\bigskip

\section{习题6.7：名义是call，支付式却是put}\label{ux4e60ux98986.7ux540dux4e49ux662fcallux652fux4ed8ux5f0fux5374ux662fput}

\textbf{书内依据：第6.2节复制、第6.3节看涨看跌平价。原题印刷248页。}\(S_0=45,u=3/2,d=2/3,R=7/6,N=2\)，执行价50。题面文字与明确的\((50-S_2)^+\)不一致，先按支付式求全部价值与持仓。

风险中性\(p^*=3/5\)。终价\(405/4,45,45,20\)，支付0、5、5、30。第一期H价值\((2/5)5/R=12/7\)，D价值\([(3/5)5+(2/5)30]/R=90/7\)，起点
\[\boxed{P_0=1296/245.}\]
初始份额\([(12/7)-(90/7)]/(135/2-30)=-52/175\)，银行\(4572/245\)。第一期H份额\(-4/45\)、银行\(54/7\)；D份额\(-1\)、银行\(300/7\)。完整核对如下。

\begingroup
\small
\setlength{\tabcolsep}{4pt}
\renewcommand{\arraystretch}{1.10}

\begin{longtable}[]{@{}
  >{\raggedleft\arraybackslash}p{(\columnwidth - 10\tabcolsep) * \real{0.1739}}
  >{\raggedright\arraybackslash}p{(\columnwidth - 10\tabcolsep) * \real{0.1304}}
  >{\raggedleft\arraybackslash}p{(\columnwidth - 10\tabcolsep) * \real{0.1739}}
  >{\raggedleft\arraybackslash}p{(\columnwidth - 10\tabcolsep) * \real{0.1739}}
  >{\raggedleft\arraybackslash}p{(\columnwidth - 10\tabcolsep) * \real{0.1739}}
  >{\raggedleft\arraybackslash}p{(\columnwidth - 10\tabcolsep) * \real{0.1739}}@{}}
\toprule\noalign{}
\begin{minipage}[b]{\linewidth}\raggedleft
期数
\end{minipage} & \begin{minipage}[b]{\linewidth}\raggedright
历史
\end{minipage} & \begin{minipage}[b]{\linewidth}\raggedleft
股价 \(S\)
\end{minipage} & \begin{minipage}[b]{\linewidth}\raggedleft
价值 \(V\)
\end{minipage} & \begin{minipage}[b]{\linewidth}\raggedleft
股数 \(\Delta\)
\end{minipage} & \begin{minipage}[b]{\linewidth}\raggedleft
银行 \(B\)
\end{minipage} \\
\midrule\noalign{}
\endhead
\bottomrule\noalign{}
\endlastfoot
0 & \(\varnothing\) & \(45\) & \(\frac{1296}{245}\) &
\(-\frac{52}{175}\) & \(\frac{4572}{245}\) \\
1 & H & \(\frac{135}{2}\) & \(\frac{12}{7}\) & \(-\frac{4}{45}\) &
\(\frac{54}{7}\) \\
1 & D & \(30\) & \(\frac{90}{7}\) & \(-1\) & \(\frac{300}{7}\) \\
\end{longtable}

\textbf{全部终端历史与支付（H为涨，D为跌）。}

\begin{longtable}[]{@{}lrrr@{}}
\toprule\noalign{}
路径 & 终价 & 到期支付 & 定价概率 \\
\midrule\noalign{}
\endhead
\bottomrule\noalign{}
\endlastfoot
HH & \(\frac{405}{4}\) & \(0\) & \(\frac{9}{25}\) \\
HD & \(45\) & \(5\) & \(\frac{6}{25}\) \\
DH & \(45\) & \(5\) & \(\frac{6}{25}\) \\
DD & \(20\) & \(30\) & \(\frac{4}{25}\) \\
\end{longtable}

\endgroup

若严格按文字``看涨''改为\((S_2-50)^+\)，不是同一合约。由平价
\[\boxed{C_0=P_0+45-50/R^2=3321/245.}\]
其股票份额在上述每个看跌份额上加1，银行减去到期付50的债券现值，即\(B_n^{C}=B_n^P-50/R^{2-n}\)。这同时给出字面解释的全部组合，不擅定作者原意。

\subsection{6.7-A：改成美式看跌}\label{aux6539ux6210ux7f8eux5f0fux770bux8dcc}

\textbf{问题。}0、1、2期都能行权，求初价和策略。

\textbf{解。}第一期H立即支付0、继续\(12/7\)，故继续；D立即支付20，大于继续\(90/7\)，故行权。起点继续值\([(3/5)(12/7)+(2/5)20]/R=\boxed{1896/245}\)，大于立即支付5。因此初始继续，第一期D就行权，H继续到期。

\subsection{6.7-B：只在最低终价支付1}\label{bux53eaux5728ux6700ux4f4eux7ec8ux4ef7ux652fux4ed81}

\textbf{问题。}求初价和第一期价值。

\textbf{解。}只有DD支付1，概率\((2/5)^2\)，故价\(\boxed{144/1225}\)。第一期H为0，D为\((2/5)/R=12/35\)。它与按整棵树``最后下跌一次''支付的数字合约不同。

\subsection{6.7-C：报价6的终值盈余}\label{cux62a5ux4ef76ux7684ux7ec8ux503cux76c8ux4f59}

\textbf{问题。}按明确看跌支付的合约报价6，理想市场中的盈余是多少？

\textbf{解。}卖出报价合约，买复制组合，初始余款\(6-1296/245=174/245\)存银行；到期所有合约风险相消，余款长为\((174/245)(7/6)^2=\boxed{29/30}\)。比较前必须先核清交易的是看跌支付，而非字面的另一份看涨。

\bigskip

\section{习题6.8：包含初价的平均价看涨}\label{ux4e60ux98986.8ux5305ux542bux521dux4ef7ux7684ux5e73ux5747ux4ef7ux770bux6da8}

\textbf{书内依据：第6.2节全路径反向复制。原题印刷248页。}\(S_0=4,u=2,d=1/2,R=5/4,N=3\)，\(A_3=(S_0+S_1+S_2+S_3)/4\)，支付\((A_3-4)^+\)。

\(p^*=1/2\)。八条路径按HHH、HHD、HDH、HDD、DHH、DHD、DDH、DDD排序，支付为\(11,5,2,1/2,1/2,0,0,0\)。例如HDD路径是\((4,8,4,2)\)，平均\(9/2\)，支付\(1/2\)；DHH路径\((4,2,4,8)\)平均也\(9/2\)。初价
\[\boxed{V_0=R^{-3}(19/8)=152/125.}\]
第二期HH、HD、DH、DD价值分别\(32/5,1,1/5,0\)，尽管HD、DH现价同为4，价值不同。第一期H、D为\(74/25,2/25\)。全部持仓与终端支付如下，每一行仍只解两个下一步方程。

\begingroup
\small
\setlength{\tabcolsep}{4pt}
\renewcommand{\arraystretch}{1.10}

\begin{longtable}[]{@{}
  >{\raggedleft\arraybackslash}p{(\columnwidth - 10\tabcolsep) * \real{0.1739}}
  >{\raggedright\arraybackslash}p{(\columnwidth - 10\tabcolsep) * \real{0.1304}}
  >{\raggedleft\arraybackslash}p{(\columnwidth - 10\tabcolsep) * \real{0.1739}}
  >{\raggedleft\arraybackslash}p{(\columnwidth - 10\tabcolsep) * \real{0.1739}}
  >{\raggedleft\arraybackslash}p{(\columnwidth - 10\tabcolsep) * \real{0.1739}}
  >{\raggedleft\arraybackslash}p{(\columnwidth - 10\tabcolsep) * \real{0.1739}}@{}}
\toprule\noalign{}
\begin{minipage}[b]{\linewidth}\raggedleft
期数
\end{minipage} & \begin{minipage}[b]{\linewidth}\raggedright
历史
\end{minipage} & \begin{minipage}[b]{\linewidth}\raggedleft
股价 \(S\)
\end{minipage} & \begin{minipage}[b]{\linewidth}\raggedleft
价值 \(V\)
\end{minipage} & \begin{minipage}[b]{\linewidth}\raggedleft
股数 \(\Delta\)
\end{minipage} & \begin{minipage}[b]{\linewidth}\raggedleft
银行 \(B\)
\end{minipage} \\
\midrule\noalign{}
\endhead
\bottomrule\noalign{}
\endlastfoot
0 & \(\varnothing\) & \(4\) & \(\frac{152}{125}\) & \(\frac{12}{25}\) &
\(-\frac{88}{125}\) \\
1 & H & \(8\) & \(\frac{74}{25}\) & \(\frac{9}{20}\) &
\(-\frac{16}{25}\) \\
1 & D & \(2\) & \(\frac{2}{25}\) & \(\frac{1}{15}\) &
\(-\frac{4}{75}\) \\
2 & HH & \(16\) & \(\frac{32}{5}\) & \(\frac{1}{4}\) &
\(\frac{12}{5}\) \\
2 & HD & \(4\) & \(1\) & \(\frac{1}{4}\) & \(0\) \\
2 & DH & \(4\) & \(\frac{1}{5}\) & \(\frac{1}{12}\) &
\(-\frac{2}{15}\) \\
2 & DD & \(1\) & \(0\) & \(0\) & \(0\) \\
\end{longtable}

\textbf{全部终端历史与支付（H为涨，D为跌）。}

\begin{longtable}[]{@{}lrrr@{}}
\toprule\noalign{}
路径 & 终价 & 到期支付 & 定价概率 \\
\midrule\noalign{}
\endhead
\bottomrule\noalign{}
\endlastfoot
HHH & \(32\) & \(11\) & \(\frac{1}{8}\) \\
HHD & \(8\) & \(5\) & \(\frac{1}{8}\) \\
HDH & \(8\) & \(2\) & \(\frac{1}{8}\) \\
HDD & \(2\) & \(\frac{1}{2}\) & \(\frac{1}{8}\) \\
DHH & \(8\) & \(\frac{1}{2}\) & \(\frac{1}{8}\) \\
DHD & \(2\) & \(0\) & \(\frac{1}{8}\) \\
DDH & \(2\) & \(0\) & \(\frac{1}{8}\) \\
DDD & \(\frac{1}{2}\) & \(0\) & \(\frac{1}{8}\) \\
\end{longtable}

\endgroup

\subsection{6.8-A：平均值不含初价}\label{aux5e73ux5747ux503cux4e0dux542bux521dux4ef7}

\textbf{问题。}改为\((S_1+S_2+S_3)/3\)，执行价仍4，求新价格。

\textbf{解。}设新平均\(A'_3\)，则\(A'_3-4=(4/3)(A_3-4)\)，由于系数正，正部支付也乘\(4/3\)。所以新价\(\boxed{608/375}\)，所有份额、现金也乘\(4/3\)。这个简化依赖初价恰好等于执行价；一般不能原样套用。

\subsection{6.8-B：与普通看涨比较}\label{bux4e0eux666eux901aux770bux6da8ux6bd4ux8f83}

\textbf{问题。}相同股票树、到期3、执行价4，普通\((S_3-4)^+\)价格是多少？

\textbf{解。}HHH支付28，三个两H路径各4，其余0，终值和40。价格\(40/(8R^3)=\boxed{64/25}\)。它大于原平均价合约，但这里通过具体树证明，不能仅凭``平均会降低波动''替代必要支付与分布核算。

\subsection{6.8-C：在D之后上涨时如何自融资换仓}\label{cux5728dux4e4bux540eux4e0aux6da8ux65f6ux5982ux4f55ux81eaux878dux8d44ux6362ux4ed3}

\textbf{问题。}检查D到DH的实际现金变化。

\textbf{解。}D节点持股\(1/15\)、银行\(-4/75\)。到DH股票涨到4，旧组合为\(4/15+(5/4)(-4/75)=1/5\)。新组合持股\(1/12\)、银行\(-2/15\)，成本\(4/12-2/15=1/5\)，无需外部资金。下一期H净值\((1/12)8+(5/4)(-2/15)=1/2\)，D净值0，恰好复制两个剩余支付。

\bigskip

\section{习题6.9：回看最低价的完整复制树}\label{ux4e60ux98986.9ux56deux770bux6700ux4f4eux4ef7ux7684ux5b8cux6574ux590dux5236ux6811}

\textbf{书内依据：第6.2节路径依赖与向后归纳。原题印刷248页。}股票与6.8相同，明确支付为\(S_3-\min_{0\le m\le3}S_m\)。题面附带的固定执行价字样不再另行扣除；按显示支付求解。

八条路径支付依次\(28,4,4,0,6,0,1,0\)，总和43，所以
\[\boxed{V_0=43/(8R^3)=344/125.}\]
例如DHH最小价2，终价8，支付6；DDH最低1、终价2，支付1。第二期HH、HD、DH、DD价值\(64/5,8/5,12/5,2/5\)；HD与DH又因历史最低价不同而不能合并。第一期H、D为\(144/25,28/25\)。

\begingroup
\small
\setlength{\tabcolsep}{4pt}
\renewcommand{\arraystretch}{1.10}

\begin{longtable}[]{@{}
  >{\raggedleft\arraybackslash}p{(\columnwidth - 10\tabcolsep) * \real{0.1739}}
  >{\raggedright\arraybackslash}p{(\columnwidth - 10\tabcolsep) * \real{0.1304}}
  >{\raggedleft\arraybackslash}p{(\columnwidth - 10\tabcolsep) * \real{0.1739}}
  >{\raggedleft\arraybackslash}p{(\columnwidth - 10\tabcolsep) * \real{0.1739}}
  >{\raggedleft\arraybackslash}p{(\columnwidth - 10\tabcolsep) * \real{0.1739}}
  >{\raggedleft\arraybackslash}p{(\columnwidth - 10\tabcolsep) * \real{0.1739}}@{}}
\toprule\noalign{}
\begin{minipage}[b]{\linewidth}\raggedleft
期数
\end{minipage} & \begin{minipage}[b]{\linewidth}\raggedright
历史
\end{minipage} & \begin{minipage}[b]{\linewidth}\raggedleft
股价 \(S\)
\end{minipage} & \begin{minipage}[b]{\linewidth}\raggedleft
价值 \(V\)
\end{minipage} & \begin{minipage}[b]{\linewidth}\raggedleft
股数 \(\Delta\)
\end{minipage} & \begin{minipage}[b]{\linewidth}\raggedleft
银行 \(B\)
\end{minipage} \\
\midrule\noalign{}
\endhead
\bottomrule\noalign{}
\endlastfoot
0 & \(\varnothing\) & \(4\) & \(\frac{344}{125}\) & \(\frac{58}{75}\) &
\(-\frac{128}{375}\) \\
1 & H & \(8\) & \(\frac{144}{25}\) & \(\frac{14}{15}\) &
\(-\frac{128}{75}\) \\
1 & D & \(2\) & \(\frac{28}{25}\) & \(\frac{2}{3}\) &
\(-\frac{16}{75}\) \\
2 & HH & \(16\) & \(\frac{64}{5}\) & \(1\) & \(-\frac{16}{5}\) \\
2 & HD & \(4\) & \(\frac{8}{5}\) & \(\frac{2}{3}\) &
\(-\frac{16}{15}\) \\
2 & DH & \(4\) & \(\frac{12}{5}\) & \(1\) & \(-\frac{8}{5}\) \\
2 & DD & \(1\) & \(\frac{2}{5}\) & \(\frac{2}{3}\) &
\(-\frac{4}{15}\) \\
\end{longtable}

\textbf{全部终端历史与支付（H为涨，D为跌）。}

\begin{longtable}[]{@{}lrrr@{}}
\toprule\noalign{}
路径 & 终价 & 到期支付 & 定价概率 \\
\midrule\noalign{}
\endhead
\bottomrule\noalign{}
\endlastfoot
HHH & \(32\) & \(28\) & \(\frac{1}{8}\) \\
HHD & \(8\) & \(4\) & \(\frac{1}{8}\) \\
HDH & \(8\) & \(4\) & \(\frac{1}{8}\) \\
HDD & \(2\) & \(0\) & \(\frac{1}{8}\) \\
DHH & \(8\) & \(6\) & \(\frac{1}{8}\) \\
DHD & \(2\) & \(0\) & \(\frac{1}{8}\) \\
DDH & \(2\) & \(1\) & \(\frac{1}{8}\) \\
DDD & \(\frac{1}{2}\) & \(0\) & \(\frac{1}{8}\) \\
\end{longtable}

\endgroup

\subsection{6.9-A：回看最高价的看跌形式}\label{aux56deux770bux6700ux9ad8ux4ef7ux7684ux770bux8dccux5f62ux5f0f}

\textbf{问题。}改为\(\max_{m\le3}S_m-S_3\)，求价格。

\textbf{解。}逐路径支付\(0,8,0,6,0,2,2,7/2\)，总和\(43/2\)。于是价\(\boxed{172/125}\)。最高价减末价和末价减最低价并非同一个合约，尽管都总为非负。

\subsection{6.9-B：整个价格区间的幅度}\label{bux6574ux4e2aux4ef7ux683cux533aux95f4ux7684ux5e45ux5ea6}

\textbf{问题。}支付为\(\max S_m-\min S_m\)，怎样无需重建整棵树定价？

\textbf{解。}逐路径\(\max S-\min S=(\max S-S_3)+(S_3-\min S)\)。两者组合即复制新支付，价格相加得\(\boxed{516/125}\)；股票和银行头寸也逐节点相加。

\subsection{\texorpdfstring{6.9-C：股票所有数值乘\(c>0\)}{6.9-C：股票所有数值乘c\textgreater0}}\label{cux80a1ux7968ux6240ux6709ux6570ux503cux4e58c0}

\textbf{问题。}回看最低价合约的价格、股数和现金怎样变？

\textbf{解。}每个支付及节点价值都乘\(c\)，复制份额分子分母同乘\(c\)故不变，银行现金乘\(c\)。初价变为\(\boxed{344c/125}\)。这里股数不变，股票投资金额仍随\(c\)放大，不能误说两者都不变。

\bigskip

\section{习题6.10：素因子个数作为支付}\label{ux4e60ux98986.10ux7d20ux56e0ux5b50ux4e2aux6570ux4f5cux4e3aux652fux4ed8}

\textbf{书内依据：第6.2节风险中性终值平均。原题印刷248页。}\(S_0=16,u=3,d=1/2,R=4/3,N=3\)，支付为终价分解中的素因子个数，重复因子重复计数。

\(p^*=(4/3-1/2)/(3-1/2)=1/3\)。若三步中有\(k\)次上涨，终价
\[S_3=16\,3^k2^{-(3-k)}=2^{k+1}3^k.\]
所以素因子数为\((k+1)+k=2k+1\)。在定价概率下\(k\sim\operatorname{Bin}(3,1/3)\)，\(Ek=1\)。平均支付3，初价
\[\boxed{V_0=3/(4/3)^3=81/64.}\]
这是素因子按重数计，而不是不同素数的种类数。为使价格不仅停留在一个期望数值，下面另列完整动态复制。

\begingroup
\small
\setlength{\tabcolsep}{4pt}
\renewcommand{\arraystretch}{1.10}

\begin{longtable}[]{@{}rlrrrr@{}}
\toprule\noalign{}
期数 & 历史 & 股价 \(S\) & 价值 \(V\) & 股数 \(\Delta\) & 银行 \(B\) \\
\midrule\noalign{}
\endhead
\bottomrule\noalign{}
\endlastfoot
0 & \(\varnothing\) & \(16\) & \(\frac{81}{64}\) & \(\frac{9}{320}\) &
\(\frac{261}{320}\) \\
1 & H & \(48\) & \(\frac{39}{16}\) & \(\frac{1}{80}\) &
\(\frac{147}{80}\) \\
1 & D & \(8\) & \(\frac{21}{16}\) & \(\frac{3}{40}\) &
\(\frac{57}{80}\) \\
2 & HH & \(144\) & \(\frac{17}{4}\) & \(\frac{1}{180}\) &
\(\frac{69}{20}\) \\
2 & HD & \(24\) & \(\frac{11}{4}\) & \(\frac{1}{30}\) &
\(\frac{39}{20}\) \\
2 & DH & \(24\) & \(\frac{11}{4}\) & \(\frac{1}{30}\) &
\(\frac{39}{20}\) \\
2 & DD & \(4\) & \(\frac{5}{4}\) & \(\frac{1}{5}\) & \(\frac{9}{20}\) \\
\end{longtable}

\textbf{全部终端历史与支付（H为涨，D为跌）。}

\begin{longtable}[]{@{}lrrr@{}}
\toprule\noalign{}
路径 & 终价 & 到期支付 & 定价概率 \\
\midrule\noalign{}
\endhead
\bottomrule\noalign{}
\endlastfoot
HHH & \(432\) & \(7\) & \(\frac{1}{27}\) \\
HHD & \(72\) & \(5\) & \(\frac{2}{27}\) \\
HDH & \(72\) & \(5\) & \(\frac{2}{27}\) \\
HDD & \(12\) & \(3\) & \(\frac{4}{27}\) \\
DHH & \(72\) & \(5\) & \(\frac{2}{27}\) \\
DHD & \(12\) & \(3\) & \(\frac{4}{27}\) \\
DDH & \(12\) & \(3\) & \(\frac{4}{27}\) \\
DDD & \(2\) & \(1\) & \(\frac{8}{27}\) \\
\end{longtable}

\endgroup

\subsection{6.10-A：只数不同素数}\label{aux53eaux6570ux4e0dux540cux7d20ux6570}

\textbf{问题。}改成不同素数的种类数，求价格。

\textbf{解。}\(k=0\)只有2一种，\(k\ge1\)有2、3两种。\(P(k=0)=(2/3)^3=8/27\)，平均支付\(2-8/27=46/27\)，折现得\(\boxed{23/32}\)。

\subsection{6.10-B：最多支付4}\label{bux6700ux591aux652fux4ed84}

\textbf{问题。}原重数支付加上上限4。

\textbf{解。}\(k=0,1,2,3\)对应支付1、3、4、4，概率\(8,12,6,1\)除以27，平均\((8+36+24+4)/27=8/3\)，价格\(\boxed{9/8}\)。封顶必须先作用于每个支付，再求期望，不能先封顶原来的平均值3。

\subsection{\texorpdfstring{6.10-C：推广到\(N\)期}{6.10-C：推广到N期}}\label{cux63a8ux5e7fux5230nux671f}

\textbf{问题。}\(S_0=2^m,m\ge N\)，每次仍乘3或\(1/2\)，求重数支付的价格。

\textbf{解。}\(k\)次H后\(S_N=2^{m-N+k}3^k\)，重数\(m-N+2k\)；\(m\ge N\)保证指数非负，终价都是正整数。由\(Ek=Np^*\)，
\[\boxed{V_0=R^{-N}(m-N+2Np^*),\quad p^*=1/3.}\]
若\(m<N\)，部分终价不是整数，必须先重新定义支付，不能延用素因子计数。

\bigskip

\section{习题6.11：三期数字期权的全树持仓}\label{ux4e60ux98986.11ux4e09ux671fux6570ux5b57ux671fux6743ux7684ux5168ux6811ux6301ux4ed3}

\textbf{书内依据：第6.2节向后定价。原题印刷248页。}\(S_0=27,u=4/3,d=2/3,R=10/9,N=3\)，到期\(S_3>27\)才付1，否则0。

\(p^*=2/3\)。\(k\)次上涨时终价\(27(4/3)^k(2/3)^{3-k}=8\,2^k\)，依次8、16、32、64。因此至少两次上涨才支付1，概率\(3(2/3)^2(1/3)+(2/3)^3=20/27\)，初价
\[\boxed{V_0=\frac{20/27}{(10/9)^3}=27/50.}\]
第一期H和D价值分别\(18/25,9/25\)，所以\(\Delta_0=(9/25)/(36-18)=1/50\)，银行\(27/50-27/50=0\)。各节点完整如下。

\begingroup
\small
\setlength{\tabcolsep}{4pt}
\renewcommand{\arraystretch}{1.10}

\begin{longtable}[]{@{}rlrrrr@{}}
\toprule\noalign{}
期数 & 历史 & 股价 \(S\) & 价值 \(V\) & 股数 \(\Delta\) & 银行 \(B\) \\
\midrule\noalign{}
\endhead
\bottomrule\noalign{}
\endlastfoot
0 & \(\varnothing\) & \(27\) & \(\frac{27}{50}\) & \(\frac{1}{50}\) &
\(0\) \\
1 & H & \(36\) & \(\frac{18}{25}\) & \(\frac{1}{80}\) &
\(\frac{27}{100}\) \\
1 & D & \(18\) & \(\frac{9}{25}\) & \(\frac{1}{20}\) &
\(-\frac{27}{50}\) \\
2 & HH & \(48\) & \(\frac{9}{10}\) & \(0\) & \(\frac{9}{10}\) \\
2 & HD & \(24\) & \(\frac{3}{5}\) & \(\frac{1}{16}\) &
\(-\frac{9}{10}\) \\
2 & DH & \(24\) & \(\frac{3}{5}\) & \(\frac{1}{16}\) &
\(-\frac{9}{10}\) \\
2 & DD & \(12\) & \(0\) & \(0\) & \(0\) \\
\end{longtable}

\textbf{全部终端历史与支付（H为涨，D为跌）。}

\begin{longtable}[]{@{}lrrr@{}}
\toprule\noalign{}
路径 & 终价 & 到期支付 & 定价概率 \\
\midrule\noalign{}
\endhead
\bottomrule\noalign{}
\endlastfoot
HHH & \(64\) & \(1\) & \(\frac{8}{27}\) \\
HHD & \(32\) & \(1\) & \(\frac{4}{27}\) \\
HDH & \(32\) & \(1\) & \(\frac{4}{27}\) \\
HDD & \(16\) & \(0\) & \(\frac{2}{27}\) \\
DHH & \(32\) & \(1\) & \(\frac{4}{27}\) \\
DHD & \(16\) & \(0\) & \(\frac{2}{27}\) \\
DDH & \(16\) & \(0\) & \(\frac{2}{27}\) \\
DDD & \(8\) & \(0\) & \(\frac{1}{27}\) \\
\end{longtable}

\endgroup

\subsection{6.11-A：门槛处的严格与非严格}\label{aux95e8ux69dbux5904ux7684ux4e25ux683cux4e0eux975eux4e25ux683c}

\textbf{问题。}分别改成\(S_3>32\)与\(S_3\ge32\)时的价。

\textbf{解。}前者只在HHH支付，概率\(8/27\)，价\(\boxed{27/125}\)；后者两次或三次上涨都付，价仍\(\boxed{27/50}\)。离散树在门槛32有正概率，不能像连续正态分布那样忽略等号。

\subsection{6.11-B：允许提前领取数字支付}\label{bux5141ux8bb8ux63d0ux524dux9886ux53d6ux6570ux5b57ux652fux4ed8}

\textbf{问题。}每期\(S_n>27\)即可行权付1，求美式价。

\textbf{解。}第二期HH现价48，立即1大于下一期折现期望\(9/10\)，故行权；HD、DH现价24，继续值\((2/3)/R=3/5\)；DD为0。第一期H立即1大于继续\([(2/3)1+(1/3)(3/5)]/R=39/50\)；D继续\([(2/3)(3/5)]/R=9/25\)。起点价值
\[\boxed{[(2/3)1+(1/3)(9/25)]/R=177/250.}\]
到期不分红普通看涨不提前行权的结论不能套到封顶数字合约。

\subsection{6.11-C：超过27时支付一股股票}\label{cux8d85ux8fc727ux65f6ux652fux4ed8ux4e00ux80a1ux80a1ux7968}

\textbf{问题。}求\(S_3\mathbf1_{\{S_3>27\}}\)。

\textbf{解。}终价64的概率\(8/27\)，32的总概率\(12/27\)。折现后的价格为
\[\boxed{R^{-3}[64(8/27)+32(12/27)]=3024/125.}\]
股数可用同一复制递推，不等于固定数字价格乘起点27，因为终价与是否达标相关。

\bigskip

\section{习题6.12：逐期观察的上敲出看跌}\label{ux4e60ux98986.12ux9010ux671fux89c2ux5bdfux7684ux4e0aux6572ux51faux770bux8dcc}

\textbf{书内依据：第6.2节路径依赖复制。原题印刷248---249页。}\(S_0=54,u=3/2,d=2/3,R=7/6,N=3,K=50\)，任何观察期价格达到上障碍70后，合约失效，未敲出时到期付\((50-S_3)^+\)。

\(p^*=3/5\)。第一期H就达到81，之后即使跌回低价也已失效。先D到36，DH到54；DHH到81触碰障碍且看跌本身也零。未敲出且正支付的路径为DHD、DDH各14，DDD为34，概率分别\(12/125,12/125,8/125\)。
\[\boxed{V_0=R^{-3}[14(24/125)+34(8/125)]=131328/42875.}\]
所有曾经敲出的节点价值与以后复制均为0；尚未敲出的节点按下一步两种结果复制。完整树保留历史，不能只按当前价格决定合约是否存活。

\begingroup
\small
\setlength{\tabcolsep}{4pt}
\renewcommand{\arraystretch}{1.10}

\begin{longtable}[]{@{}
  >{\raggedleft\arraybackslash}p{(\columnwidth - 10\tabcolsep) * \real{0.1739}}
  >{\raggedright\arraybackslash}p{(\columnwidth - 10\tabcolsep) * \real{0.1304}}
  >{\raggedleft\arraybackslash}p{(\columnwidth - 10\tabcolsep) * \real{0.1739}}
  >{\raggedleft\arraybackslash}p{(\columnwidth - 10\tabcolsep) * \real{0.1739}}
  >{\raggedleft\arraybackslash}p{(\columnwidth - 10\tabcolsep) * \real{0.1739}}
  >{\raggedleft\arraybackslash}p{(\columnwidth - 10\tabcolsep) * \real{0.1739}}@{}}
\toprule\noalign{}
\begin{minipage}[b]{\linewidth}\raggedleft
期数
\end{minipage} & \begin{minipage}[b]{\linewidth}\raggedright
历史
\end{minipage} & \begin{minipage}[b]{\linewidth}\raggedleft
股价 \(S\)
\end{minipage} & \begin{minipage}[b]{\linewidth}\raggedleft
价值 \(V\)
\end{minipage} & \begin{minipage}[b]{\linewidth}\raggedleft
股数 \(\Delta\)
\end{minipage} & \begin{minipage}[b]{\linewidth}\raggedleft
银行 \(B\)
\end{minipage} \\
\midrule\noalign{}
\endhead
\bottomrule\noalign{}
\endlastfoot
0 & \(\varnothing\) & \(54\) & \(\frac{131328}{42875}\) &
\(-\frac{1216}{6125}\) & \(\frac{590976}{42875}\) \\
1 & H & \(81\) & \(0\) & \(0\) & \(0\) \\
1 & D & \(36\) & \(\frac{10944}{1225}\) & \(-\frac{82}{175}\) &
\(\frac{31608}{1225}\) \\
2 & HH & \(\frac{243}{2}\) & \(0\) & \(0\) & \(0\) \\
2 & HD & \(54\) & \(0\) & \(0\) & \(0\) \\
2 & DH & \(54\) & \(\frac{24}{5}\) & \(-\frac{14}{45}\) &
\(\frac{108}{5}\) \\
2 & DD & \(24\) & \(\frac{132}{7}\) & \(-1\) & \(\frac{300}{7}\) \\
\end{longtable}

\textbf{全部终端历史与支付（H为涨，D为跌）。}

\begin{longtable}[]{@{}lrrr@{}}
\toprule\noalign{}
路径 & 终价 & 到期支付 & 定价概率 \\
\midrule\noalign{}
\endhead
\bottomrule\noalign{}
\endlastfoot
HHH & \(\frac{729}{4}\) & \(0\) & \(\frac{27}{125}\) \\
HHD & \(81\) & \(0\) & \(\frac{18}{125}\) \\
HDH & \(81\) & \(0\) & \(\frac{18}{125}\) \\
HDD & \(36\) & \(0\) & \(\frac{12}{125}\) \\
DHH & \(81\) & \(0\) & \(\frac{18}{125}\) \\
DHD & \(36\) & \(14\) & \(\frac{12}{125}\) \\
DDH & \(36\) & \(14\) & \(\frac{12}{125}\) \\
DDD & \(16\) & \(34\) & \(\frac{8}{125}\) \\
\end{longtable}

\endgroup

\subsection{6.12-A：配对的上敲入看跌}\label{aux914dux5bf9ux7684ux4e0aux6572ux5165ux770bux8dcc}

\textbf{问题。}只有曾达到70时才有效的同执行价看跌值多少？

\textbf{解。}每条路径不是敲入就是敲出，两份相加等于普通看跌。普通看跌价\(167616/42875\)，减原题得\(\boxed{5184/6125}\)。也可直接看唯一被敲出且本来正支付的HDD：概率\((3/5)(2/5)^2=12/125\)、支付14，折现即同值。

\subsection{6.12-B：只在到期检查障碍}\label{bux53eaux5728ux5230ux671fux68c0ux67e5ux969cux788d}

\textbf{问题。}若只判断\(S_3\ge70\)，价格如何？

\textbf{解。}普通看跌有正支付时\(S_3<50<70\)，所以不会被这种到期障碍删除。价就是普通看跌\(\boxed{167616/42875}\)。中途曾超过70又跌下来的路径原题不付，这个新合约却付，观察频率改变了支付。

\subsection{6.12-C：敲出后到期返还2}\label{cux6572ux51faux540eux5230ux671fux8fd4ux8fd82}

\textbf{问题。}仍逐期观察，敲出后统一到期付2，求价格。

\textbf{解。}敲出事件为第一步H，或第一步D随后HH。概率\(3/5+(2/5)(3/5)^2=93/125\)。增加的现值为\(2(93/125)/R^3\)，加原值得\(\boxed{171504/42875}\)。若返还在敲出当时支付，折现时刻不同，不能沿用这个计算。

\bigskip

\section{习题6.13：障碍合约与一个同价格的非障碍支付}\label{ux4e60ux98986.13ux969cux788dux5408ux7ea6ux4e0eux4e00ux4e2aux540cux4ef7ux683cux7684ux975eux969cux788dux652fux4ed8}

\textbf{书内依据：第6.2节风险中性计价与路径树。原题印刷249页。}\(S_0=32,u=2,d=1/2,R=5/4,N=4\)，上障碍100，看跌执行价50，逐期敲出。

\subsection{（a）同价的终价函数怎样构造}\label{aux540cux4ef7ux7684ux7ec8ux4ef7ux51fdux6570ux600eux6837ux6784ux9020}

\(p^*=1/2\)，16条历史等概率。普通看跌的正支付路径中，只有HHDD曾经越过100：价格轨迹32、64、128、64、32，本来付18，却被删去。因此敲出价等于普通看跌价减去``路径HHDD上支付18''的现值。

构造一个仅依赖终价的支付：除\(S_4=512\)时付\(-18\)之外，其他终价仍付\((50-S_4)^+\)。终价512仅由HHHH产生，和HHDD的概率同为\(1/16\)，故将负18从一条路径挪到另一条路径不改变定价期望。这就是（a）要求的同价非障碍支付。

它们只同初价，\textbf{不是逐路径相等}：HHDD在新支付下给18而原敲出给0；HHHH新支付\(-18\)而原敲出0。

\subsection{（b）实际价格}\label{bux5b9eux9645ux4ef7ux683c}

普通正支付总和为：全D终价2，付48；一次H的四条各付42；两次H的六条各付18。删掉HHDD的18后总和\(48+4\cdot42+5\cdot18=306\)。因此
\[\boxed{V_0=306/(16R^4)=4896/625.}\]
完整复制仍按原障碍支付而不是同价替代支付计算。

\begingroup
\small
\setlength{\tabcolsep}{4pt}
\renewcommand{\arraystretch}{1.10}

\begin{longtable}[]{@{}
  >{\raggedleft\arraybackslash}p{(\columnwidth - 10\tabcolsep) * \real{0.1739}}
  >{\raggedright\arraybackslash}p{(\columnwidth - 10\tabcolsep) * \real{0.1304}}
  >{\raggedleft\arraybackslash}p{(\columnwidth - 10\tabcolsep) * \real{0.1739}}
  >{\raggedleft\arraybackslash}p{(\columnwidth - 10\tabcolsep) * \real{0.1739}}
  >{\raggedleft\arraybackslash}p{(\columnwidth - 10\tabcolsep) * \real{0.1739}}
  >{\raggedleft\arraybackslash}p{(\columnwidth - 10\tabcolsep) * \real{0.1739}}@{}}
\toprule\noalign{}
\begin{minipage}[b]{\linewidth}\raggedleft
期数
\end{minipage} & \begin{minipage}[b]{\linewidth}\raggedright
历史
\end{minipage} & \begin{minipage}[b]{\linewidth}\raggedleft
股价 \(S\)
\end{minipage} & \begin{minipage}[b]{\linewidth}\raggedleft
价值 \(V\)
\end{minipage} & \begin{minipage}[b]{\linewidth}\raggedleft
股数 \(\Delta\)
\end{minipage} & \begin{minipage}[b]{\linewidth}\raggedleft
银行 \(B\)
\end{minipage} \\
\midrule\noalign{}
\endhead
\bottomrule\noalign{}
\endlastfoot
0 & \(\varnothing\) & \(32\) & \(\frac{4896}{625}\) & \(-\frac{1}{5}\) &
\(\frac{8896}{625}\) \\
1 & H & \(64\) & \(\frac{624}{125}\) & \(-\frac{13}{100}\) &
\(\frac{1664}{125}\) \\
1 & D & \(16\) & \(\frac{1824}{125}\) & \(-\frac{12}{25}\) &
\(\frac{2784}{125}\) \\
2 & HH & \(128\) & \(0\) & \(0\) & \(0\) \\
2 & HD & \(32\) & \(\frac{312}{25}\) & \(-\frac{7}{20}\) &
\(\frac{592}{25}\) \\
2 & DH & \(32\) & \(\frac{312}{25}\) & \(-\frac{7}{20}\) &
\(\frac{592}{25}\) \\
2 & DD & \(8\) & \(24\) & \(-1\) & \(32\) \\
3 & HHH & \(256\) & \(0\) & \(0\) & \(0\) \\
3 & HHD & \(64\) & \(0\) & \(0\) & \(0\) \\
3 & HDH & \(64\) & \(\frac{36}{5}\) & \(-\frac{3}{16}\) &
\(\frac{96}{5}\) \\
3 & HDD & \(16\) & \(24\) & \(-1\) & \(40\) \\
3 & DHH & \(64\) & \(\frac{36}{5}\) & \(-\frac{3}{16}\) &
\(\frac{96}{5}\) \\
3 & DHD & \(16\) & \(24\) & \(-1\) & \(40\) \\
3 & DDH & \(16\) & \(24\) & \(-1\) & \(40\) \\
3 & DDD & \(4\) & \(36\) & \(-1\) & \(40\) \\
\end{longtable}

\textbf{全部终端历史与支付（H为涨，D为跌）。}

\begin{longtable}[]{@{}lrrr@{}}
\toprule\noalign{}
路径 & 终价 & 到期支付 & 定价概率 \\
\midrule\noalign{}
\endhead
\bottomrule\noalign{}
\endlastfoot
HHHH & \(512\) & \(0\) & \(\frac{1}{16}\) \\
HHHD & \(128\) & \(0\) & \(\frac{1}{16}\) \\
HHDH & \(128\) & \(0\) & \(\frac{1}{16}\) \\
HHDD & \(32\) & \(0\) & \(\frac{1}{16}\) \\
HDHH & \(128\) & \(0\) & \(\frac{1}{16}\) \\
HDHD & \(32\) & \(18\) & \(\frac{1}{16}\) \\
HDDH & \(32\) & \(18\) & \(\frac{1}{16}\) \\
HDDD & \(8\) & \(42\) & \(\frac{1}{16}\) \\
DHHH & \(128\) & \(0\) & \(\frac{1}{16}\) \\
DHHD & \(32\) & \(18\) & \(\frac{1}{16}\) \\
DHDH & \(32\) & \(18\) & \(\frac{1}{16}\) \\
DHDD & \(8\) & \(42\) & \(\frac{1}{16}\) \\
DDHH & \(32\) & \(18\) & \(\frac{1}{16}\) \\
DDHD & \(8\) & \(42\) & \(\frac{1}{16}\) \\
DDDH & \(8\) & \(42\) & \(\frac{1}{16}\) \\
DDDD & \(2\) & \(48\) & \(\frac{1}{16}\) \\
\end{longtable}

\endgroup

\subsection{6.13-A：普通看跌与敲入部分}\label{aux666eux901aux770bux8dccux4e0eux6572ux5165ux90e8ux5206}

\textbf{问题。}分别求普通、敲入两种价格。

\textbf{解。}普通总支付和324，价\(\boxed{5184/625}\)；敲入只含HHDD的18，价\(\boxed{288/625}\)。两者之差恢复原题。配对恒等式是逐路径成立，而（a）构造只有初价相等。

\subsection{6.13-B：上涨权重不再一半}\label{bux4e0aux6da8ux6743ux91cdux4e0dux518dux4e00ux534a}

\textbf{问题。}若定价上涨概率\(p\in(0,1)\)，将HHDD上删除的18转移成HHHH上的负支付，应付多少？

\textbf{解。}两路径概率分别\(p^2q^2,p^4\)，要期望相同，负支付幅度\(b\)须满足\(bp^4=18p^2q^2\)。所以\(\boxed{b=18(q/p)^2}\)。等概率时的18不是任意概率下都成立。

\subsection{6.13-C：相同初价不保证中途同价}\label{cux76f8ux540cux521dux4ef7ux4e0dux4fddux8bc1ux4e2dux9014ux540cux4ef7}

\textbf{问题。}在第三期HHH节点比较原障碍合约与（a）的替代合约。

\textbf{解。}原合约早已在第二期HH越障碍，价值0；替代合约下一步HHHH付\(-18\)，HHHD终价128付0，故节点值\((-18+0)/(2R)=\boxed{-36/5}\)。两者中途价格不同，不能只凭相同初始期望把复制组合互换。

\bigskip

\section{习题6.14：看涨复制与相同执行价看跌}\label{ux4e60ux98986.14ux770bux6da8ux590dux5236ux4e0eux76f8ux540cux6267ux884cux4ef7ux770bux8dcc}

\textbf{书内依据：第6.2节复制、第6.3节定理6.6平价。原题印刷249页。}\(S_0=64,u=2,d=1/2,R=5/4,N=3,K=125\)。

\subsection{（a）看涨价格}\label{aux770bux6da8ux4ef7ux683c}

\(p^*=1/2\)。HHH终价512，支付387；两次H的三条路径终价128，各付3；其余零。故
\[\boxed{C_0=(387+9)/(8R^3)=3168/125.}\] \#\# （b）所有时刻持仓

第二期HH价值156，HD和DH为\(6/5\)，DD为0；第一期H、D为\(1572/25,12/25\)。起点份额\(13/20\)，银行\(-2032/125\)。各行含相应的现金负债，逐节点无需再加外部资金。

\begingroup
\small
\setlength{\tabcolsep}{4pt}
\renewcommand{\arraystretch}{1.10}

\begin{longtable}[]{@{}
  >{\raggedleft\arraybackslash}p{(\columnwidth - 10\tabcolsep) * \real{0.1739}}
  >{\raggedright\arraybackslash}p{(\columnwidth - 10\tabcolsep) * \real{0.1304}}
  >{\raggedleft\arraybackslash}p{(\columnwidth - 10\tabcolsep) * \real{0.1739}}
  >{\raggedleft\arraybackslash}p{(\columnwidth - 10\tabcolsep) * \real{0.1739}}
  >{\raggedleft\arraybackslash}p{(\columnwidth - 10\tabcolsep) * \real{0.1739}}
  >{\raggedleft\arraybackslash}p{(\columnwidth - 10\tabcolsep) * \real{0.1739}}@{}}
\toprule\noalign{}
\begin{minipage}[b]{\linewidth}\raggedleft
期数
\end{minipage} & \begin{minipage}[b]{\linewidth}\raggedright
历史
\end{minipage} & \begin{minipage}[b]{\linewidth}\raggedleft
股价 \(S\)
\end{minipage} & \begin{minipage}[b]{\linewidth}\raggedleft
价值 \(V\)
\end{minipage} & \begin{minipage}[b]{\linewidth}\raggedleft
股数 \(\Delta\)
\end{minipage} & \begin{minipage}[b]{\linewidth}\raggedleft
银行 \(B\)
\end{minipage} \\
\midrule\noalign{}
\endhead
\bottomrule\noalign{}
\endlastfoot
0 & \(\varnothing\) & \(64\) & \(\frac{3168}{125}\) & \(\frac{13}{20}\)
& \(-\frac{2032}{125}\) \\
1 & H & \(128\) & \(\frac{1572}{25}\) & \(\frac{129}{160}\) &
\(-\frac{1008}{25}\) \\
1 & D & \(32\) & \(\frac{12}{25}\) & \(\frac{1}{40}\) &
\(-\frac{8}{25}\) \\
2 & HH & \(256\) & \(156\) & \(1\) & \(-100\) \\
2 & HD & \(64\) & \(\frac{6}{5}\) & \(\frac{1}{32}\) &
\(-\frac{4}{5}\) \\
2 & DH & \(64\) & \(\frac{6}{5}\) & \(\frac{1}{32}\) &
\(-\frac{4}{5}\) \\
2 & DD & \(16\) & \(0\) & \(0\) & \(0\) \\
\end{longtable}

\textbf{全部终端历史与支付（H为涨，D为跌）。}

\begin{longtable}[]{@{}lrrr@{}}
\toprule\noalign{}
路径 & 终价 & 到期支付 & 定价概率 \\
\midrule\noalign{}
\endhead
\bottomrule\noalign{}
\endlastfoot
HHH & \(512\) & \(387\) & \(\frac{1}{8}\) \\
HHD & \(128\) & \(3\) & \(\frac{1}{8}\) \\
HDH & \(128\) & \(3\) & \(\frac{1}{8}\) \\
HDD & \(32\) & \(0\) & \(\frac{1}{8}\) \\
DHH & \(128\) & \(3\) & \(\frac{1}{8}\) \\
DHD & \(32\) & \(0\) & \(\frac{1}{8}\) \\
DDH & \(32\) & \(0\) & \(\frac{1}{8}\) \\
DDD & \(8\) & \(0\) & \(\frac{1}{8}\) \\
\end{longtable}

\endgroup

\subsection{（c）看跌价格}\label{cux770bux8dccux4ef7ux683c}

执行价现值\(125/R^3=64=S_0\)，所以平价\(C_0-P_0=S_0-125/R^3=0\)，得到\(\boxed{P_0=3168/125}\)。这是当前特定参数使两价相等，不是看涨看跌一般等价。

\subsection{6.14-A：数字支付}\label{aux6570ux5b57ux652fux4ed8}

\textbf{问题。}\(S_3>125\)时付1。

\textbf{解。}至少两次H的四条路径支付，概率\(1/2\)，所以\(\boxed{V_0=(1/2)/R^3=32/125}\)。四条支付大小相同，不需要再乘终价128或512。

\subsection{6.14-B：看涨最多支付100}\label{bux770bux6da8ux6700ux591aux652fux4ed8100}

\textbf{问题。}求\(\min\{(S_3-125)^+,100\}\)的价。

\textbf{解。}HHH支付从387降为100，两H的三条仍各3，总和109。所以价\(\boxed{872/125}\)。也可将普通看涨减掉执行价225的看涨，逐路径相同。

\subsection{6.14-C：一般何时看涨与看跌同价}\label{cux4e00ux822cux4f55ux65f6ux770bux6da8ux4e0eux770bux8dccux540cux4ef7}

\textbf{问题。}无分红\(N\)期树，求相同执行价\(K\)的两价相等的充要条件。

\textbf{解。}逐路径\(C_N-P_N=S_N-K\)，复制价格线性给\(C_0-P_0=S_0-K/R^N\)，故\(\boxed{K=S_0R^N}\)当且仅当两价相同。结论不依赖上涨、下跌具体因子，只要基础模型无套利。

\bigskip

\section{习题6.15：欧式看跌完整三期复制}\label{ux4e60ux98986.15ux6b27ux5f0fux770bux8dccux5b8cux6574ux4e09ux671fux590dux5236}

\textbf{书内依据：第6.2节自融资复制。原题印刷249页。}\(S_0=27,u=4/3,d=2/3,R=10/9,N=3,K=30\)，与6.11相同股票树。

\(p^*=2/3\)。终价按H次数为8、16、32、64；看跌分别22、14、0、0。第二期HH、HD、DH、DD价值为\(0,21/5,21/5,15\)；第一期H、D为\(63/50,351/50\)。所以
\[\boxed{P_0=[(2/3)(63/50)+(1/3)(351/50)]/R=1431/500.}\]
初始份额为\(\Delta_0=-8/25\)，银行\(B_0=5751/500\)。第一期H份额\(-7/40\)、银行\(189/25\)；D份额\(-9/10\)、银行\(1161/50\)。第二期HD、DH份额\(-7/8\)、银行\(126/5\)；DD份额\(-1\)、银行27。

\begingroup
\small
\setlength{\tabcolsep}{4pt}
\renewcommand{\arraystretch}{1.10}

\begin{longtable}[]{@{}
  >{\raggedleft\arraybackslash}p{(\columnwidth - 10\tabcolsep) * \real{0.1739}}
  >{\raggedright\arraybackslash}p{(\columnwidth - 10\tabcolsep) * \real{0.1304}}
  >{\raggedleft\arraybackslash}p{(\columnwidth - 10\tabcolsep) * \real{0.1739}}
  >{\raggedleft\arraybackslash}p{(\columnwidth - 10\tabcolsep) * \real{0.1739}}
  >{\raggedleft\arraybackslash}p{(\columnwidth - 10\tabcolsep) * \real{0.1739}}
  >{\raggedleft\arraybackslash}p{(\columnwidth - 10\tabcolsep) * \real{0.1739}}@{}}
\toprule\noalign{}
\begin{minipage}[b]{\linewidth}\raggedleft
期数
\end{minipage} & \begin{minipage}[b]{\linewidth}\raggedright
历史
\end{minipage} & \begin{minipage}[b]{\linewidth}\raggedleft
股价 \(S\)
\end{minipage} & \begin{minipage}[b]{\linewidth}\raggedleft
价值 \(V\)
\end{minipage} & \begin{minipage}[b]{\linewidth}\raggedleft
股数 \(\Delta\)
\end{minipage} & \begin{minipage}[b]{\linewidth}\raggedleft
银行 \(B\)
\end{minipage} \\
\midrule\noalign{}
\endhead
\bottomrule\noalign{}
\endlastfoot
0 & \(\varnothing\) & \(27\) & \(\frac{1431}{500}\) & \(-\frac{8}{25}\)
& \(\frac{5751}{500}\) \\
1 & H & \(36\) & \(\frac{63}{50}\) & \(-\frac{7}{40}\) &
\(\frac{189}{25}\) \\
1 & D & \(18\) & \(\frac{351}{50}\) & \(-\frac{9}{10}\) &
\(\frac{1161}{50}\) \\
2 & HH & \(48\) & \(0\) & \(0\) & \(0\) \\
2 & HD & \(24\) & \(\frac{21}{5}\) & \(-\frac{7}{8}\) &
\(\frac{126}{5}\) \\
2 & DH & \(24\) & \(\frac{21}{5}\) & \(-\frac{7}{8}\) &
\(\frac{126}{5}\) \\
2 & DD & \(12\) & \(15\) & \(-1\) & \(27\) \\
\end{longtable}

\textbf{全部终端历史与支付（H为涨，D为跌）。}

\begin{longtable}[]{@{}lrrr@{}}
\toprule\noalign{}
路径 & 终价 & 到期支付 & 定价概率 \\
\midrule\noalign{}
\endhead
\bottomrule\noalign{}
\endlastfoot
HHH & \(64\) & \(0\) & \(\frac{8}{27}\) \\
HHD & \(32\) & \(0\) & \(\frac{4}{27}\) \\
HDH & \(32\) & \(0\) & \(\frac{4}{27}\) \\
HDD & \(16\) & \(14\) & \(\frac{2}{27}\) \\
DHH & \(32\) & \(0\) & \(\frac{4}{27}\) \\
DHD & \(16\) & \(14\) & \(\frac{2}{27}\) \\
DDH & \(16\) & \(14\) & \(\frac{2}{27}\) \\
DDD & \(8\) & \(22\) & \(\frac{1}{27}\) \\
\end{longtable}

\endgroup

\subsection{6.15-A：同时买看涨看跌}\label{aux540cux65f6ux4e70ux770bux6da8ux770bux8dcc}

\textbf{问题。}求同执行价看涨及欧式跨式\(|S_3-30|\)的价值。

\textbf{解。}执行价现值\(30/R^3=2187/100\)，由平价看涨\(1431/500+27-2187/100=\boxed{999/125}\)。跨式逐路径为两者之和，所以\(\boxed{5427/500}\)。

\subsection{6.15-B：把实际上涨概率误当定价概率}\label{bux628aux5b9eux9645ux4e0aux6da8ux6982ux7387ux8befux5f53ux5b9aux4ef7ux6982ux7387}

\textbf{问题。}即使预测实际上涨概率为\(1/2\)，可否用该权重折现支付定价？

\textbf{解。}等权终端支付和64，均值8，折现为\(\boxed{729/125}\)，不同于复制价\(1431/500\)。因为\((u+d)/2=1\)不等于银行因子\(10/9\)，该权重不使股票折现均值不变。复制成本由两个状态财富方程决定，不随主观概率改变。

\subsection{6.15-C：初始组合买好就不再动}\label{cux521dux59cbux7ec4ux5408ux4e70ux597dux5c31ux4e0dux518dux52a8}

\textbf{问题。}一直持有初始\(\Delta=-8/25,B=5751/500\)，是否仍可复制？

\textbf{解。}在HHH终价64，静态组合终值
\[-\frac8{25}64+\frac{5751}{500}\left(\frac{10}{9}\right)^3
=-\frac{512}{25}+\frac{1278}{81}<0,\]
但期权支付0。原初始组合只保证下一期两个\textbf{继续价值}，之后必须按树重新配置，不是静态复制全部终端。

\bigskip

\section{习题6.16：美式看跌与美式看涨}\label{ux4e60ux98986.16ux7f8eux5f0fux770bux8dccux4e0eux7f8eux5f0fux770bux6da8}

\textbf{书内依据：第6.4节式(6.18)---(6.19)和定理6.7。原题印刷249页。}使用6.15的股票树，0至3期都能行权，执行价30。

\subsection{（a）美式看跌逐步比较}\label{aux7f8eux5f0fux770bux8dccux9010ux6b65ux6bd4ux8f83}

第二期HH现价48，支付0，继续也0；HD、DH现价24，立即得6，大于继续\(21/5\)；DD现价12，立即18，大于继续15。因此第二期节点值为\(0,6,6,18\)。

第一期H继续值\([(2/3)0+(1/3)6]/R=9/5\)，立即0，故继续；D立即\(30-18=12\)，大于继续\([(2/3)6+(1/3)18]/R=9\)，故行权。起点继续价值
\[\boxed{P_0^A=[(2/3)(9/5)+(1/3)12]/R=117/25}\]
大于立即支付3。所以实际最优策略为起点继续，第一期D行权；第一期H之后第二期HD行权，HH价值0，立即或到期都可。下表也列了实际策略已退出后仍可数学定义的其他节点，不能把那些后续支付再次累计。

\begingroup
\small
\setlength{\tabcolsep}{4pt}
\renewcommand{\arraystretch}{1.10}

\begin{longtable}[]{@{}rlrrrrl@{}}
\toprule\noalign{}
期数 & 历史 & 股价 \(S\) & 立即 \(g\) & 继续 \(C\) & 最优 \(V\) &
选择 \\
\midrule\noalign{}
\endhead
\bottomrule\noalign{}
\endlastfoot
0 & \(\varnothing\) & \(27\) & \(3\) & \(\frac{117}{25}\) &
\(\frac{117}{25}\) & 继续 \\
1 & H & \(36\) & \(0\) & \(\frac{9}{5}\) & \(\frac{9}{5}\) & 继续 \\
1 & D & \(18\) & \(12\) & \(9\) & \(12\) & 行权 \\
2 & HH & \(48\) & \(0\) & \(0\) & \(0\) & 两者均可 \\
2 & HD & \(24\) & \(6\) & \(\frac{21}{5}\) & \(6\) & 行权 \\
2 & DH & \(24\) & \(6\) & \(\frac{21}{5}\) & \(6\) & 行权 \\
2 & DD & \(12\) & \(18\) & \(15\) & \(18\) & 行权 \\
\end{longtable}

\textbf{全部终端历史与支付（H为涨，D为跌）。}

\begin{longtable}[]{@{}lrrr@{}}
\toprule\noalign{}
路径 & 终价 & 到期支付 & 定价概率 \\
\midrule\noalign{}
\endhead
\bottomrule\noalign{}
\endlastfoot
HHH & \(64\) & \(0\) & \(\frac{8}{27}\) \\
HHD & \(32\) & \(0\) & \(\frac{4}{27}\) \\
HDH & \(32\) & \(0\) & \(\frac{4}{27}\) \\
HDD & \(16\) & \(14\) & \(\frac{2}{27}\) \\
DHH & \(32\) & \(0\) & \(\frac{4}{27}\) \\
DHD & \(16\) & \(14\) & \(\frac{2}{27}\) \\
DDH & \(16\) & \(14\) & \(\frac{2}{27}\) \\
DDD & \(8\) & \(22\) & \(\frac{1}{27}\) \\
\end{longtable}

\endgroup

\subsection{（b）美式看涨为什么不需提前行权}\label{bux7f8eux5f0fux770bux6da8ux4e3aux4ec0ux4e48ux4e0dux9700ux63d0ux524dux884cux6743}

无分红且\(R\ge1\)。在任一第\(n\)期节点，剩\(m\)期，若坚持到期，欧式值为\(R^{-m}E^*(S_N-K)^+\)。凸函数\((x-K)^+\)的期望至少为均值的正部，且\(E^*S_N=R^mS_n\)，所以
\[R^{-m}E^*(S_N-K)^+\ge(S_n-K/R^m)^+\ge(S_n-K)^+.\]
因此等待到期的价值不低于立即行权。每个中间节点都成立，逆推可知美式价值等于欧式，\(\boxed{C_0^A=999/125}\)。下面给实际各节点比较，和一般证明互相核对。

\begingroup
\small
\setlength{\tabcolsep}{4pt}
\renewcommand{\arraystretch}{1.10}

\begin{longtable}[]{@{}rlrrrrl@{}}
\toprule\noalign{}
期数 & 历史 & 股价 \(S\) & 立即 \(g\) & 继续 \(C\) & 最优 \(V\) &
选择 \\
\midrule\noalign{}
\endhead
\bottomrule\noalign{}
\endlastfoot
0 & \(\varnothing\) & \(27\) & \(0\) & \(\frac{999}{125}\) &
\(\frac{999}{125}\) & 继续 \\
1 & H & \(36\) & \(6\) & \(\frac{324}{25}\) & \(\frac{324}{25}\) &
继续 \\
1 & D & \(18\) & \(0\) & \(\frac{18}{25}\) & \(\frac{18}{25}\) & 继续 \\
2 & HH & \(48\) & \(18\) & \(21\) & \(21\) & 继续 \\
2 & HD & \(24\) & \(0\) & \(\frac{6}{5}\) & \(\frac{6}{5}\) & 继续 \\
2 & DH & \(24\) & \(0\) & \(\frac{6}{5}\) & \(\frac{6}{5}\) & 继续 \\
2 & DD & \(12\) & \(0\) & \(0\) & \(0\) & 两者均可 \\
\end{longtable}

\textbf{全部终端历史与支付（H为涨，D为跌）。}

\begin{longtable}[]{@{}lrrr@{}}
\toprule\noalign{}
路径 & 终价 & 到期支付 & 定价概率 \\
\midrule\noalign{}
\endhead
\bottomrule\noalign{}
\endlastfoot
HHH & \(64\) & \(34\) & \(\frac{8}{27}\) \\
HHD & \(32\) & \(2\) & \(\frac{4}{27}\) \\
HDH & \(32\) & \(2\) & \(\frac{4}{27}\) \\
HDD & \(16\) & \(0\) & \(\frac{2}{27}\) \\
DHH & \(32\) & \(2\) & \(\frac{4}{27}\) \\
DHD & \(16\) & \(0\) & \(\frac{2}{27}\) \\
DDH & \(16\) & \(0\) & \(\frac{2}{27}\) \\
DDD & \(8\) & \(0\) & \(\frac{1}{27}\) \\
\end{longtable}

\endgroup

\subsection{6.16-A：只在第2或3期行权}\label{aux53eaux5728ux7b2c2ux62163ux671fux884cux6743}

\textbf{问题。}求这份Bermudan看跌价值。

\textbf{解。}第二期节点值仍\(0,6,6,18\)。第一期虽然D的立即值12更高，却不允许行权，只能使用继续9；H仍\(9/5\)。故起点\([(2/3)(9/5)+(1/3)9]/R=\boxed{189/50}\)。它介于欧式\(1431/500\)和美式\(117/25\)之间。

\subsection{6.16-B：负利率反例}\label{bux8d1fux5229ux7387ux53cdux4f8b}

\textbf{问题。}股票100下一期110或90，\(R=19/20\)，看涨执行价80。比较现在行权和等一期。

\textbf{解。}\(d=.9<R=.95<u=1.1\)仍无套利，风险中性上涨概率1/4。两个支付30、10，欧式价值\([(1/4)30+(3/4)10]/(.95)=300/19<20\)。现在行权获得20，所以美式价\(\boxed{20}\)。原``不提前行权''证明确实用到了\(R\ge1\)。

\subsection{6.16-C：随机决定是否行权能否更好}\label{cux968fux673aux51b3ux5b9aux662fux5426ux884cux6743ux80fdux5426ux66f4ux597d}

\textbf{问题。}某节点以概率\(\alpha\)行权、其余继续，是否可超过取最大？

\textbf{解。}若两个条件价值分别\(g,C\)，混合值\(\alpha g+(1-\alpha)C\le\max(g,C)\)。从最后一期逐层用此不等式，任何随机策略都不优于每节点取最大。相等时随机化可以同样最优，不能提高价值。

\bigskip

\section{习题6.17：美式跨式不能随意拆成两个美式期权}\label{ux4e60ux98986.17ux7f8eux5f0fux8de8ux5f0fux4e0dux80fdux968fux610fux62c6ux6210ux4e24ux4e2aux7f8eux5f0fux671fux6743}

\textbf{书内依据：第6.4节美式向后归纳。原题印刷249页。}仍用\(S_0=27,u=4/3,d=2/3,R=10/9,N=3\)，行权支付\(|S_n-30|\)，只允许领取一次。

\subsection{全树比较与初价}\label{ux5168ux6811ux6bd4ux8f83ux4e0eux521dux4ef7}

第二期HH的现价为\(27(4/3)^2=48\)，立即支付\(48-30=18\)。下一期64、32支付34、2，继续\([(2/3)34+(1/3)2]/R=21\)，故\textbf{继续，节点值21}。HD、DH现价24，立即6，大于继续\(27/5\)；DD现价12，立即18，大于继续15。第二期最优值为\(21,6,6,18\)。

第一期H现价36，立即6，继续\([(2/3)21+(1/3)6]/R=72/5\)，故继续；D立即12大于继续9，故行权。起点继续值
\[\boxed{V_0=[(2/3)(72/5)+(1/3)12]/R=306/25}\]
大于立即3。最优策略是D在第一期行权；H分支走到HD在第二期行权；HH继续到期。注意HH的\textbf{继续21}不能错写成``立即21''；完整表以实际支付计算为准。

\begingroup
\small
\setlength{\tabcolsep}{4pt}
\renewcommand{\arraystretch}{1.10}

\begin{longtable}[]{@{}rlrrrrl@{}}
\toprule\noalign{}
期数 & 历史 & 股价 \(S\) & 立即 \(g\) & 继续 \(C\) & 最优 \(V\) &
选择 \\
\midrule\noalign{}
\endhead
\bottomrule\noalign{}
\endlastfoot
0 & \(\varnothing\) & \(27\) & \(3\) & \(\frac{306}{25}\) &
\(\frac{306}{25}\) & 继续 \\
1 & H & \(36\) & \(6\) & \(\frac{72}{5}\) & \(\frac{72}{5}\) & 继续 \\
1 & D & \(18\) & \(12\) & \(9\) & \(12\) & 行权 \\
2 & HH & \(48\) & \(18\) & \(21\) & \(21\) & 继续 \\
2 & HD & \(24\) & \(6\) & \(\frac{27}{5}\) & \(6\) & 行权 \\
2 & DH & \(24\) & \(6\) & \(\frac{27}{5}\) & \(6\) & 行权 \\
2 & DD & \(12\) & \(18\) & \(15\) & \(18\) & 行权 \\
\end{longtable}

\textbf{全部终端历史与支付（H为涨，D为跌）。}

\begin{longtable}[]{@{}lrrr@{}}
\toprule\noalign{}
路径 & 终价 & 到期支付 & 定价概率 \\
\midrule\noalign{}
\endhead
\bottomrule\noalign{}
\endlastfoot
HHH & \(64\) & \(34\) & \(\frac{8}{27}\) \\
HHD & \(32\) & \(2\) & \(\frac{4}{27}\) \\
HDH & \(32\) & \(2\) & \(\frac{4}{27}\) \\
HDD & \(16\) & \(14\) & \(\frac{2}{27}\) \\
DHH & \(32\) & \(2\) & \(\frac{4}{27}\) \\
DHD & \(16\) & \(14\) & \(\frac{2}{27}\) \\
DDH & \(16\) & \(14\) & \(\frac{2}{27}\) \\
DDD & \(8\) & \(22\) & \(\frac{1}{27}\) \\
\end{longtable}

\endgroup

\subsection{两种美式权利为何不是简单相加}\label{ux4e24ux79cdux7f8eux5f0fux6743ux5229ux4e3aux4f55ux4e0dux662fux7b80ux5355ux76f8ux52a0}

逐状态\(|S_n-K|=(S_n-K)^++(K-S_n)^+\)，但单份跨式必须在同一时刻同时结束两部分。分开的美式看涨与看跌可以各选自己的停止时刻，因此分开购买价值不小于跨式。本题
\[C_0^A+P_0^A=999/125+117/25=1584/125,\]
而跨式\(306/25=1530/125\)，严格少\(\boxed{54/125}\)。不能把``支付相加''与``各自独立优化停止时间''交换顺序。

\subsection{6.17-A：欧式跨式与提前行权溢价}\label{aux6b27ux5f0fux8de8ux5f0fux4e0eux63d0ux524dux884cux6743ux6ea2ux4ef7}

\textbf{问题。}求两者及差额。

\textbf{解。}欧式跨式可直接相加欧式看涨看跌，值\(5427/500\)。美式\(306/25=6120/500\)，溢价\(\boxed{693/500}\)。欧式两部分的停止日期相同，线性相加没有障碍。

\subsection{6.17-B：相对分开两份美式，损失来自哪些路径}\label{bux76f8ux5bf9ux5206ux5f00ux4e24ux4efdux7f8eux5f0fux635fux5931ux6765ux81eaux54eaux4e9bux8defux5f84}

\textbf{问题。}解释差额\(54/125\)。

\textbf{解。}第一期D时跨式随看跌一起结束，放弃仍有价值\(18/25\)的欧式看涨；此事件折现权重\((1/3)/R=3/10\)，损失\(27/125\)。第二期HD也随看跌结束，放弃看涨\(6/5\)，事件权重\((2/3)(1/3)/R^2=9/50\)，损失另\(27/125\)。相加\(54/125\)；两事件互斥，不重复计账。

\subsection{6.17-C：何时合并价恰好等于分开价}\label{cux4f55ux65f6ux5408ux5e76ux4ef7ux6070ux597dux7b49ux4e8eux5206ux5f00ux4ef7}

\textbf{问题。}有限树上两份美式支付，证明相加价格等号成立的条件。

\textbf{解。}树上策略有限，最优策略都能达到。若存在一个同时对两份支付最优的停止策略，按此策略领取两者即可达到各自最优价之和，故相等。反过来，若合并价等于二者之和，取达到合并最优的停止策略；它在每一份上的价值不超过对应单独最优值，两个非负差额之和为0，因此每个差额都0，此策略同时最优。故等号当且仅当存在共同最优停止策略。

\bigskip

\section{习题6.18：封顶美式看涨可以提前行权}\label{ux4e60ux98986.18ux5c01ux9876ux7f8eux5f0fux770bux6da8ux53efux4ee5ux63d0ux524dux884cux6743}

\textbf{书内依据：第6.4节动态规划。原题印刷249页。}\(S_0=16,u=3,d=1/2,R=4/3,N=3\)，\(p^*=1/3\)，每期可领取\(g(S_n)=\min\{(S_n-10)^+,60\}\)。

\subsection{从最后一期倒推}\label{ux4eceux6700ux540eux4e00ux671fux5012ux63a8}

第二期HH现价144，立即已到上限60，继续45，故领取60；HD、DH现价24，立即14小于继续16；DD现价4，立即0小于继续\(1/2\)。

第一期H现价48，立即38，大于继续\([(1/3)60+(2/3)16]/R=23\)，故行权；D现价8，立即0小于继续\([(1/3)16+(2/3)(1/2)]/R=17/4\)。起点继续值
\[\boxed{V_0=[(1/3)38+(2/3)(17/4)]/R=93/8}\]
大于立即6。实际策略：第一期H就领取38，第一期D分支到期领取；虽表内HH建议领取60，但从初始最优策略已经不会走到仍持有合约的HH。

\begingroup
\small
\setlength{\tabcolsep}{4pt}
\renewcommand{\arraystretch}{1.10}

\begin{longtable}[]{@{}rlrrrrl@{}}
\toprule\noalign{}
期数 & 历史 & 股价 \(S\) & 立即 \(g\) & 继续 \(C\) & 最优 \(V\) &
选择 \\
\midrule\noalign{}
\endhead
\bottomrule\noalign{}
\endlastfoot
0 & \(\varnothing\) & \(16\) & \(6\) & \(\frac{93}{8}\) &
\(\frac{93}{8}\) & 继续 \\
1 & H & \(48\) & \(38\) & \(23\) & \(38\) & 行权 \\
1 & D & \(8\) & \(0\) & \(\frac{17}{4}\) & \(\frac{17}{4}\) & 继续 \\
2 & HH & \(144\) & \(60\) & \(45\) & \(60\) & 行权 \\
2 & HD & \(24\) & \(14\) & \(16\) & \(16\) & 继续 \\
2 & DH & \(24\) & \(14\) & \(16\) & \(16\) & 继续 \\
2 & DD & \(4\) & \(0\) & \(\frac{1}{2}\) & \(\frac{1}{2}\) & 继续 \\
\end{longtable}

\textbf{全部终端历史与支付（H为涨，D为跌）。}

\begin{longtable}[]{@{}lrrr@{}}
\toprule\noalign{}
路径 & 终价 & 到期支付 & 定价概率 \\
\midrule\noalign{}
\endhead
\bottomrule\noalign{}
\endlastfoot
HHH & \(432\) & \(60\) & \(\frac{1}{27}\) \\
HHD & \(72\) & \(60\) & \(\frac{2}{27}\) \\
HDH & \(72\) & \(60\) & \(\frac{2}{27}\) \\
HDD & \(12\) & \(2\) & \(\frac{4}{27}\) \\
DHH & \(72\) & \(60\) & \(\frac{2}{27}\) \\
DHD & \(12\) & \(2\) & \(\frac{4}{27}\) \\
DDH & \(12\) & \(2\) & \(\frac{4}{27}\) \\
DDD & \(2\) & \(0\) & \(\frac{8}{27}\) \\
\end{longtable}

\endgroup

\subsection{6.18-A：如果只能到期领取}\label{aux5982ux679cux53eaux80fdux5230ux671fux9886ux53d6}

\textbf{问题。}求欧式价与提前行权溢价。

\textbf{解。}终端按H次数0、1、2、3的支付为0、2、60、60，概率\(8,12,6,1\)除以27。平均支付\([24+360+60]/27=148/9\)，折现后\(\boxed{111/16}\)。美式\(93/8=186/16\)，溢价\(\boxed{75/16}\)。

\subsection{6.18-B：达到上限时为什么立即领取}\label{bux8fbeux5230ux4e0aux9650ux65f6ux4e3aux4ec0ux4e48ux7acbux5373ux9886ux53d6}

\textbf{问题。}任何未来支付都至多\(L>0\)，当前立即支付恰为\(L\)，且\(R>1\)，证明现在行权最优。

\textbf{解。}若至少等一期，即使以后确定得到最大支付\(L\)，折现回现在也至多\(L/R<L\)。随机等待更久只进一步折现。故严格应现在领取。封顶破坏了普通看涨未来无限上涨的好处，不能套无分红普通call的定理。

\subsection{6.18-C：按最优停止直接求价格和时间}\label{cux6309ux6700ux4f18ux505cux6b62ux76f4ux63a5ux6c42ux4ef7ux683cux548cux65f6ux95f4}

\textbf{问题。}不再逆推，按实际策略直接复算，并求停止期数均值。

\textbf{解。}第一期H概率1/3，得38；否则概率2/3，到第3期。故\(E^*\tau=(1/3)1+(2/3)3=\boxed{7/3}\)。在D分支，DHH付60、概率2/27，DHD与DDH各付2、各概率4/27，DDD付0。因此
\[\frac{38/3}{R}+\frac{60(2/27)+2(8/27)}{R^3}=\boxed{93/8}.\]
必须每笔按真实停止时刻折现，不能把38也折三期。

\bigskip

\section{习题6.19：美式回看最高价合约}\label{ux4e60ux98986.19ux7f8eux5f0fux56deux770bux6700ux9ad8ux4ef7ux5408ux7ea6}

\textbf{书内依据：第6.4节美式递推与历史状态扩充。原题印刷249页。}\(S_0=8,u=2,d=1/2,R=5/4,N=3,p^*=1/2\)，行权支付\(\max_{m\le n}S_m-S_n\)。

第二期HH当前32、历史最高32，立即0，继续\(32/5\)；HD当前8、最高16，立即8大于继续\(24/5\)；DH当前8、最高8，立即0小于继续\(8/5\)；DD当前2、最高8，立即6大于继续\(22/5\)。所以第二期值\(32/5,8,8/5,6\)。

第一期H的继续值\((32/5+8)/(2R)=144/25\)，立即0；D立即\(8-4=4\)，大于继续\((8/5+6)/(2R)=76/25\)。起点
\[\boxed{V_0=(144/25+4)/(2R)=488/125.}\]
最优实际策略为第一期D行权，或先H再D在第二期行权，否则到期。下面每个历史状态都单列。

\begingroup
\small
\setlength{\tabcolsep}{4pt}
\renewcommand{\arraystretch}{1.10}

\begin{longtable}[]{@{}rlrrrrl@{}}
\toprule\noalign{}
期数 & 历史 & 股价 \(S\) & 立即 \(g\) & 继续 \(C\) & 最优 \(V\) &
选择 \\
\midrule\noalign{}
\endhead
\bottomrule\noalign{}
\endlastfoot
0 & \(\varnothing\) & \(8\) & \(0\) & \(\frac{488}{125}\) &
\(\frac{488}{125}\) & 继续 \\
1 & H & \(16\) & \(0\) & \(\frac{144}{25}\) & \(\frac{144}{25}\) &
继续 \\
1 & D & \(4\) & \(4\) & \(\frac{76}{25}\) & \(4\) & 行权 \\
2 & HH & \(32\) & \(0\) & \(\frac{32}{5}\) & \(\frac{32}{5}\) & 继续 \\
2 & HD & \(8\) & \(8\) & \(\frac{24}{5}\) & \(8\) & 行权 \\
2 & DH & \(8\) & \(0\) & \(\frac{8}{5}\) & \(\frac{8}{5}\) & 继续 \\
2 & DD & \(2\) & \(6\) & \(\frac{22}{5}\) & \(6\) & 行权 \\
\end{longtable}

\textbf{全部终端历史与支付（H为涨，D为跌）。}

\begin{longtable}[]{@{}lrrr@{}}
\toprule\noalign{}
路径 & 终价 & 到期支付 & 定价概率 \\
\midrule\noalign{}
\endhead
\bottomrule\noalign{}
\endlastfoot
HHH & \(64\) & \(0\) & \(\frac{1}{8}\) \\
HHD & \(16\) & \(16\) & \(\frac{1}{8}\) \\
HDH & \(16\) & \(0\) & \(\frac{1}{8}\) \\
HDD & \(4\) & \(12\) & \(\frac{1}{8}\) \\
DHH & \(16\) & \(0\) & \(\frac{1}{8}\) \\
DHD & \(4\) & \(4\) & \(\frac{1}{8}\) \\
DDH & \(4\) & \(4\) & \(\frac{1}{8}\) \\
DDD & \(1\) & \(7\) & \(\frac{1}{8}\) \\
\end{longtable}

\endgroup

\subsection{6.19-A：固定到期的对应合约}\label{aux56faux5b9aux5230ux671fux7684ux5bf9ux5e94ux5408ux7ea6}

\textbf{问题。}求欧式价及美式增加的价值。

\textbf{解。}终端回看支付为0、16、0、12、0、4、4、7，总和43，故欧式价\(43/(8R^3)=\boxed{344/125}\)。美式比它多\(\boxed{144/125}\)。增加的是选择不同领取时刻的权利，不是改变终端公式。

\subsection{6.19-B：同股价却不同价值}\label{bux540cux80a1ux4ef7ux5374ux4e0dux540cux4ef7ux503c}

\textbf{问题。}比较第二期HD和DH。

\textbf{解。}二者都为股价8，但HD曾到16，立即付8，最优价值8；DH历史最高8，立即0，最优继续价值\(8/5\)。因此股票现价一个变量不足以作为本合约完整状态，至少还要记历史最高价。

\subsection{6.19-C：支付比例调整}\label{cux652fux4ed8ux6bd4ux4f8bux8c03ux6574}

\textbf{问题。}改成\(c(\max_{m\le n}S_m-S_n)\)，\(c>0\)。

\textbf{解。}末期支付乘\(c\)。若下一层价值都乘\(c\)，继续值也乘\(c\)，且\(\max(cg,cC)=c\max(g,C)\)，由归纳所有节点同倍放大，最优比较不变。因此初价\(\boxed{488c/125}\)，行权策略不变。\(c=0\)时所有策略同样好，应另说并列。

\bigskip

\section{习题6.20：平均价与现价支付的两种完整解读}\label{ux4e60ux98986.20ux5e73ux5747ux4ef7ux4e0eux73b0ux4ef7ux652fux4ed8ux7684ux4e24ux79cdux5b8cux6574ux89e3ux8bfb}

\textbf{书内依据：第6.4节美式动态规划、第6.2节路径树。原题印刷250页。}\(S_0=4,u=2,d=1/2,R=5/4,N=3\)。题面称Asian并定义\(A_n=(S_0+\cdots+S_n)/(n+1)\)，随后支付式却写\((S_n-4)^+\)。本稿不静默替换符号，分别求解。

\subsection{\texorpdfstring{解读一：Asian含义，按平均价\((A_n-4)^+\)行权}{解读一：Asian含义，按平均价(A\_n-4)\^{}+行权}}\label{ux89e3ux8bfbux4e00asianux542bux4e49ux6309ux5e73ux5747ux4ef7a_n-4ux884cux6743}

末期支付与6.8相同。第二期HH平均\(28/3\)，立即\(16/3\)小于继续\(32/5\)；HD平均\(16/3\)，立即\(4/3\)大于继续1；DH平均\(10/3\)，立即0小于继续\(1/5\)；DD立即与继续都是0。因此第二期最优值\(32/5,4/3,1/5,0\)。

第一期H平均6、立即2，继续\((32/5+4/3)/(2R)=232/75\)；D立即0、继续\(2/25\)。起点立即0，继续
\[\boxed{V_0^{\rm Asian}=(232/75+2/25)/(2R)=476/375.}\]
实际提前行权只发生在第二期HD；其余继续或零值并列。

\begingroup
\small
\setlength{\tabcolsep}{4pt}
\renewcommand{\arraystretch}{1.10}

\begin{longtable}[]{@{}rlrrrrl@{}}
\toprule\noalign{}
期数 & 历史 & 股价 \(S\) & 立即 \(g\) & 继续 \(C\) & 最优 \(V\) &
选择 \\
\midrule\noalign{}
\endhead
\bottomrule\noalign{}
\endlastfoot
0 & \(\varnothing\) & \(4\) & \(0\) & \(\frac{476}{375}\) &
\(\frac{476}{375}\) & 继续 \\
1 & H & \(8\) & \(2\) & \(\frac{232}{75}\) & \(\frac{232}{75}\) &
继续 \\
1 & D & \(2\) & \(0\) & \(\frac{2}{25}\) & \(\frac{2}{25}\) & 继续 \\
2 & HH & \(16\) & \(\frac{16}{3}\) & \(\frac{32}{5}\) & \(\frac{32}{5}\)
& 继续 \\
2 & HD & \(4\) & \(\frac{4}{3}\) & \(1\) & \(\frac{4}{3}\) & 行权 \\
2 & DH & \(4\) & \(0\) & \(\frac{1}{5}\) & \(\frac{1}{5}\) & 继续 \\
2 & DD & \(1\) & \(0\) & \(0\) & \(0\) & 两者均可 \\
\end{longtable}

\textbf{全部终端历史与支付（H为涨，D为跌）。}

\begin{longtable}[]{@{}lrrr@{}}
\toprule\noalign{}
路径 & 终价 & 到期支付 & 定价概率 \\
\midrule\noalign{}
\endhead
\bottomrule\noalign{}
\endlastfoot
HHH & \(32\) & \(11\) & \(\frac{1}{8}\) \\
HHD & \(8\) & \(5\) & \(\frac{1}{8}\) \\
HDH & \(8\) & \(2\) & \(\frac{1}{8}\) \\
HDD & \(2\) & \(\frac{1}{2}\) & \(\frac{1}{8}\) \\
DHH & \(8\) & \(\frac{1}{2}\) & \(\frac{1}{8}\) \\
DHD & \(2\) & \(0\) & \(\frac{1}{8}\) \\
DDH & \(2\) & \(0\) & \(\frac{1}{8}\) \\
DDD & \(\frac{1}{2}\) & \(0\) & \(\frac{1}{8}\) \\
\end{longtable}

\endgroup

\subsection{\texorpdfstring{解读二：严格按印出的\((S_n-4)^+\)}{解读二：严格按印出的(S\_n-4)\^{}+}}\label{ux89e3ux8bfbux4e8cux4e25ux683cux6309ux5370ux51faux7684s_n-4}

这是普通美式看涨，不再是平均价合约。无分红且\(R>1\)，由6.16的证明，到期策略不劣于提前行权，因此价值等于6.8-B的欧式普通看涨，
\[\boxed{V_0^{\rm spot}=64/25.}\]
下面给第二种合约的完整比较表。两种答案不等，是支付定义不同，不是两种算法算同一个问题得到不一致。

\begingroup
\small
\setlength{\tabcolsep}{4pt}
\renewcommand{\arraystretch}{1.10}

\begin{longtable}[]{@{}rlrrrrl@{}}
\toprule\noalign{}
期数 & 历史 & 股价 \(S\) & 立即 \(g\) & 继续 \(C\) & 最优 \(V\) &
选择 \\
\midrule\noalign{}
\endhead
\bottomrule\noalign{}
\endlastfoot
0 & \(\varnothing\) & \(4\) & \(0\) & \(\frac{64}{25}\) &
\(\frac{64}{25}\) & 继续 \\
1 & H & \(8\) & \(4\) & \(\frac{144}{25}\) & \(\frac{144}{25}\) &
继续 \\
1 & D & \(2\) & \(0\) & \(\frac{16}{25}\) & \(\frac{16}{25}\) & 继续 \\
2 & HH & \(16\) & \(12\) & \(\frac{64}{5}\) & \(\frac{64}{5}\) & 继续 \\
2 & HD & \(4\) & \(0\) & \(\frac{8}{5}\) & \(\frac{8}{5}\) & 继续 \\
2 & DH & \(4\) & \(0\) & \(\frac{8}{5}\) & \(\frac{8}{5}\) & 继续 \\
2 & DD & \(1\) & \(0\) & \(0\) & \(0\) & 两者均可 \\
\end{longtable}

\textbf{全部终端历史与支付（H为涨，D为跌）。}

\begin{longtable}[]{@{}lrrr@{}}
\toprule\noalign{}
路径 & 终价 & 到期支付 & 定价概率 \\
\midrule\noalign{}
\endhead
\bottomrule\noalign{}
\endlastfoot
HHH & \(32\) & \(28\) & \(\frac{1}{8}\) \\
HHD & \(8\) & \(4\) & \(\frac{1}{8}\) \\
HDH & \(8\) & \(4\) & \(\frac{1}{8}\) \\
HDD & \(2\) & \(0\) & \(\frac{1}{8}\) \\
DHH & \(8\) & \(4\) & \(\frac{1}{8}\) \\
DHD & \(2\) & \(0\) & \(\frac{1}{8}\) \\
DDH & \(2\) & \(0\) & \(\frac{1}{8}\) \\
DDD & \(\frac{1}{2}\) & \(0\) & \(\frac{1}{8}\) \\
\end{longtable}

\endgroup

\subsection{6.20-A：Asian提前行权溢价}\label{aasianux63d0ux524dux884cux6743ux6ea2ux4ef7}

\textbf{问题。}比较解读一的美式与6.8的欧式。

\textbf{解。}欧式\(152/125=456/375\)，美式\(476/375\)，多\(\boxed{4/75}\)。只在HD节点提升价值：立即\(4/3\)减原继续1为\(1/3\)，该历史定价概率1/4，折现两期给\((1/3)(1/4)/R^2=4/75\)，直接复算同值。

\subsection{6.20-B：相同现价，平均价不同}\label{bux76f8ux540cux73b0ux4ef7ux5e73ux5747ux4ef7ux4e0dux540c}

\textbf{问题。}比较HD和DH的立即值与最优值。

\textbf{解。}两条现价都4；HD均价\(16/3\)，立即\(4/3\)，最优也\(4/3\)；DH均价\(10/3\)，立即0，最优继续\(1/5\)。由此反证若只保存股票现价便会遗漏决定行权的信息。

\subsection{6.20-C：平均不含初价，不能统一乘常数}\label{cux5e73ux5747ux4e0dux542bux521dux4ef7ux4e0dux80fdux7edfux4e00ux4e58ux5e38ux6570}

\textbf{问题。}第\(n\ge1\)期平均改为\(A'_n=(S_1+\cdots+S_n)/n\)，0期不行权，求美式价。

\textbf{解。}由于\(S_0=4\)，有\((A'_n-4)^+=(n+1)(A_n-4)^+/n\)。放大系数随\(n\)变，不能把原美式初价统一乘\(4/3\)。末期是原来的\(4/3\)倍；逆推后第二期HH、HD、DH、DD值为\(128/15,2,4/15,0\)，第一期H、D为\(316/75,8/75\)（H立即4仍小于继续\(316/75\)）。所以
\[\boxed{V_0=(316/75+8/75)/(2R)=216/125.}\]
每个日期的平均定义都要先固定，再逐层优化。

\bigskip

\section{习题6.21：直接验证Black--Scholes方程}\label{ux4e60ux98986.21ux76f4ux63a5ux9a8cux8bc1blackscholesux65b9ux7a0b}

\textbf{书内依据：第6.5---6.6节无分红Black--Scholes方程。原题印刷250页。}验证\(V(t,s)=as+be^{rt}\)满足
\[V_t+rsV_s+\frac12\sigma^2s^2V_{ss}-rV=0.\]

逐项求偏导：\(V_t=rbe^{rt},V_s=a,V_{ss}=0\)。代入左边为
\[rbe^{rt}+rsa+0-r(as+be^{rt})=0,\]
所以确为解。它的金融含义也完全明确：始终持\(a\)股、银行初始放\(b\)，\(t\)时银行值\(be^{rt}\)，不用调整持仓。这里\(t\)是向前的日历时间；不能把银行项误写成\(be^{-rt}\)。

\subsection{6.21-A：指定终端线性支付}\label{aux6307ux5b9aux7ec8ux7aefux7ebfux6027ux652fux4ed8}

\textbf{问题。}到\(T\)支付\(aS_T+b\)，求\(0\le t\le T\)时价值。

\textbf{解。}持\(a\)股，加一份到\(T\)正好变成\(b\)的债券，故
\[\boxed{V(t,s)=as+be^{-r(T-t)}.}\]
\(t=T\)恰为支付，且这是原题解中的银行初值改成\(be^{-rT}\)。终端常数与0时银行存款不能混为一个\(b\)。

\subsection{6.21-B：幂次终值}\label{bux5e42ux6b21ux7ec8ux503c}

\textbf{问题。}股票始终正，终端支付\(S_T^m\)，\(m\)为实数，找一个方程解。

\textbf{解。}尝试\(V(t,s)=s^me^{k(T-t)}\)，则\(V_t=-kV,V_s=mV/s,V_{ss}=m(m-1)V/s^2\)。代入得\([-k+rm+\sigma^2m(m-1)/2-r]V=0\)，故
\[\boxed{V=s^m\exp\{[(m-1)r+\tfrac12\sigma^2m(m-1)](T-t)\}.}\]
\(m=0\)回到零息债券，\(m=1\)回到股票本身。对数正态分布所有实数幂矩有限，风险中性积分也给此解。

\subsection{6.21-C：把折现提前吸收进去}\label{cux628aux6298ux73b0ux63d0ux524dux5438ux6536ux8fdbux53bb}

\textbf{问题。}定义\(U(t,s)=e^{-rt}V(t,s)\)，求它的方程。

\textbf{解。}反写\(V=e^{rt}U\)，所以\(V_t=e^{rt}(rU+U_t),V_s=e^{rt}U_s,V_{ss}=e^{rt}U_{ss}\)。代入原式，\(rU\)与\(-rU\)抵消，除去正因子得
\[\boxed{U_t+rsU_s+\tfrac12\sigma^2s^2U_{ss}=0.}\]
变换去掉零阶项，但终端条件也变成\(U(T,s)=e^{-rT}g(s)\)，不能漏改。

\bigskip

\section{习题6.22：二元看涨的连续模型价格}\label{ux4e60ux98986.22ux4e8cux5143ux770bux6da8ux7684ux8fdeux7eedux6a21ux578bux4ef7ux683c}

\textbf{书内依据：第6.6节对数正态终值、贴现期望与正态积分。原题印刷250页。}到期股价超过执行价才付1；一般公式后，代入执行价等于现价、剩余3个月、\(r=0,\sigma=0.3\)。

\subsection{第一步：把到期达标写成正态事件}\label{ux7b2cux4e00ux6b65ux628aux5230ux671fux8fbeux6807ux5199ux6210ux6b63ux6001ux4e8bux4ef6}

风险中性终值为\(S_\tau=S_0e^{(r-\sigma^2/2)\tau+\sigma\sqrt\tau Z}\)。对\(S_\tau>K\)两边取对数并移项，得到
\[Z>\frac{\log(K/S_0)-(r-\sigma^2/2)\tau}{\sigma\sqrt\tau}=-d_2.\]
标准正态对称性给\(P^*(Z>-d_2)=\Phi(d_2)\)，故现金数字合约现值
\[\boxed{D=e^{-r\tau}\Phi(d_2).}\]
\(\sigma\sqrt\tau>0\)时等号事件概率0，严格与非严格门槛相同；退化波动或立即到期需直接看确定支付。

\subsection{第二步：代入而不混淆三个月与三年}\label{ux7b2cux4e8cux6b65ux4ee3ux5165ux800cux4e0dux6df7ux6dc6ux4e09ux4e2aux6708ux4e0eux4e09ux5e74}

\(\tau=1/4,K=S_0,r=0\)，因此\(d_2=-(\sigma/2)\sqrt\tau=-0.075\)，所以
\[\boxed{D=\Phi(-0.075)\approx0.470107355947.}\]
不等于1/2：风险中性股票均值虽然保持不变，对数均值却有\(-\sigma^2\tau/2\)修正。

\subsection{6.22-A：达标后付股票而不是1元}\label{aux8fbeux6807ux540eux4ed8ux80a1ux7968ux800cux4e0dux662f1ux5143}

\textbf{问题。}求\(S_\tau\mathbf1_{\{S_\tau>K\}}\)的价值。

\textbf{解。}令\(v=\sigma\sqrt\tau\)，价值为\(S_0\int_{-d_2}^{\infty}e^{-v^2/2+vz}\varphi(z)dz\)。指数完成平方给\(e^{-v^2/2+vz}\varphi(z)=\varphi(z-v)\)，换元后积分下限\(-d_2-v=-d_1\)，所以\(\boxed{S_0\Phi(d_1)}\)。它与原数字合约不同；两者相减适当倍数恰好得到普通call公式。

\subsection{6.22-B：窄看涨价差趋于数字合约}\label{bux7a84ux770bux6da8ux4ef7ux5deeux8d8bux4e8eux6570ux5b57ux5408ux7ea6}

\textbf{问题。}证明\([C(K)-C(K+\varepsilon)]/\varepsilon\)在\(\varepsilon\downarrow0\)时趋于原价。

\textbf{解。}逐终值\(y\)，\([(y-K)^+-(y-K-\varepsilon)^+]/\varepsilon\)位于0与1之间，除\(y=K\)外趋于\(\mathbf1_{\{y>K\}}\)。终值连续，门槛点零概率，有界收敛给价差极限\(e^{-r\tau}\Phi(d_2)\)。因此\(-\partial C/\partial K=D\)，不能凭形式求导却漏掉随\(K\)变化的两个正态项。

\subsection{6.22-C：数字合约的对冲敏感度}\label{cux6570ux5b57ux5408ux7ea6ux7684ux5bf9ux51b2ux654fux611fux5ea6}

\textbf{问题。}求对现价的导数，观察平值且期限趋零时的行为。

\textbf{解。}\(\partial d_2/\partial S_0=1/(S_0\sigma\sqrt\tau)\)，链式法则给
\[\boxed{\Delta_D=\frac{e^{-r\tau}\varphi(d_2)}{S_0\sigma\sqrt\tau}.}\]
平值时\(d_2\to0\)，故\(\Delta_D\sim1/(S_0\sigma\sqrt{2\pi\tau})\)，趋无穷。这是数学模型中到期附近跳跃支付的敏感性，不意味着现实中可以无风险实施无限快速对冲。

\bigskip

\section{习题6.23：长期看涨合约的数值价格}\label{ux4e60ux98986.23ux957fux671fux770bux6da8ux5408ux7ea6ux7684ux6570ux503cux4ef7ux683c}

\textbf{书内依据：第6.6节Black--Scholes公式。原题印刷250页。}本题指定历史模型输入\(S_0=74.625,K=100,\tau=1.646\)年，\(r=0.05,\sigma=0.375\)，求欧式看涨价。这些是教材数值，不是当前市场报价。

\subsection{列出中间量再代公式}\label{ux5217ux51faux4e2dux95f4ux91cfux518dux4ee3ux516cux5f0f}

\[\sigma\sqrt\tau\approx0.481111993199,\qquad
Ke^{-r\tau}\approx92.099561856183.\]
\[d_1=\frac{\log(74.625/100)+(0.05+0.375^2/2)1.646}{0.375\sqrt{1.646}}
\approx-0.196753023440,\]
\[d_2=d_1-0.375\sqrt{1.646}\approx-0.677865016639.\]
查正态积分或数值积分得到\(\Phi(d_1)\approx0.422010407150,\Phi(d_2)\approx0.248928641414\)，故
\[\boxed{C=74.625\Phi(d_1)-92.099561856183\Phi(d_2)
\approx8.566307825875.}\]
现价低于执行价不意味着看涨价值0，因为未来仍有达标可能。数值取舍以给定模型为准，不用今天的数据更新这道历史练习。

\subsection{6.23-A：由平价求看跌}\label{aux7531ux5e73ux4ef7ux6c42ux770bux8dcc}

\textbf{问题。}同参数看跌是多少？

\textbf{解。}\(P=C+Ke^{-r\tau}-S_0\)，所以\(\boxed{P\approx26.040869682058}\)。这也等于\(Ke^{-r\tau}\Phi(-d_2)-S_0\Phi(-d_1)\)，正态对称性使两式完全一致。

\subsection{6.23-B：风险中性概率不是实际预测概率}\label{bux98ceux9669ux4e2dux6027ux6982ux7387ux4e0dux662fux5b9eux9645ux9884ux6d4bux6982ux7387}

\textbf{问题。}另作假设，真实漂移为\(a=0.08\)，其他相同，真实达标概率怎样？

\textbf{解。}真实模型的对数漂移改成\(a-\sigma^2/2\)，所以
\[\boxed{P(S_\tau>K)=\Phi\left(\frac{\log(S_0/K)+(a-\sigma^2/2)\tau}{\sigma\sqrt\tau}\right)\approx0.282568624743.}\]
定价权重下的概率则为\(\Phi(d_2)\approx0.248928641414\)。这里的0.08只是变式假设，不是对该公司真实收益的调查或预测。

\subsection{6.23-C：波动率增加一点的影响}\label{cux6ce2ux52a8ux7387ux589eux52a0ux4e00ux70b9ux7684ux5f71ux54cd}

\textbf{问题。}求Vega并估计\(\sigma\)增加0.01后的价格变化。

\textbf{解。}直接展开\(d_1,d_2\)可验证\(S_0\varphi(d_1)=Ke^{-r\tau}\varphi(d_2)\)。对call公式求\(\sigma\)导数，两个共同项抵消，只留下
\[\boxed{\frac{\partial C}{\partial\sigma}=S_0\varphi(d_1)\sqrt\tau\approx37.46305373267.}\]
因此一阶近似增加\(0.01\)时价格增加\(\boxed{0.3746305373}\)。这只是导数的一阶估计，不冒称变化0.01后精确价格差。

\bigskip

\section{习题6.24：一个看涨报价合理，不代表看跌也合理}\label{ux4e60ux98986.24ux4e00ux4e2aux770bux6da8ux62a5ux4ef7ux5408ux7406ux4e0dux4ee3ux8868ux770bux8dccux4e5fux5408ux7406}

\textbf{书内依据：第6.6节连续模型，第6.3节看涨看跌平价。原题印刷250页。}给定\(S_0=36.83,K=33,\tau=0.227,r=0.01,\sigma=0.15\)，教材看涨报价区间3.9---4.0、看跌报价0.4，比较模型值。

\subsection{（a）看涨价格}\label{aux770bux6da8ux4ef7ux683c-1}

\[\sigma\sqrt\tau\approx0.071466775497,\quad
Ke^{-r\tau}\approx32.925174958555,\]
\[d_1\approx1.60394698365,\qquad d_2\approx1.53248020816.\]
代公式得\(\boxed{C\approx3.96711451102}\)，确在给定3.9---4.0之间。

\subsection{（b）看跌价格与报价的差别}\label{bux770bux8dccux4ef7ux683cux4e0eux62a5ux4ef7ux7684ux5deeux522b}

用平价而不再重复正态积分：
\[\boxed{P=C+32.925174958555-36.83\approx0.062289469574.}\]
显著低于0.4。即使不指定波动率，只要求看涨价格落在给定区间，平价能允许的看跌上界也只有\(4+32.925174958555-36.83=0.095174958555\)。所以在无分红、同一到期和执行价的理想模型内，这些看涨、看跌与股票价格不同时相容。原区间下端3.9还略低于看涨内在下界\(S_0-PV(K)\)，不可把其对应负看跌值当作合法价格。

\subsection{6.24-A：把假设写明再构造组合}\label{aux628aux5047ux8bbeux5199ux660eux518dux6784ux9020ux7ec4ux5408}

\textbf{问题。}若能按4买看涨、按0.4卖看跌，并可按36.83卖股票，构造理想套利。

\textbf{解。}卖看跌和股票，买看涨及到期付33的债券。初始收入
\[0.4+36.83-4-32.925174958555=\boxed{0.304825041445}.\]
终值\(-(33-S_T)^+-S_T+(S_T-33)^++33=0\)，所以正初始盈余可存银行。这里特意增加``确能按这些方向成交''的假设；仅有一个历史报价数字不自动保证现实无成本可成交。

\subsection{6.24-B：保持看跌0.4，需要什么看涨价}\label{bux4fddux6301ux770bux8dcc0.4ux9700ux8981ux4ec0ux4e48ux770bux6da8ux4ef7}

\textbf{问题。}股票、债券和看跌不变，只改看涨。

\textbf{解。}平价要求\(C=P+S_0-PV(K)\)，所以\(\boxed{C=4.304825041445}\)。它超出给定区间，修改波动率不能改变这条线性恒等式。

\subsection{6.24-C：假设存在连续股息率会怎样}\label{cux5047ux8bbeux5b58ux5728ux8fdeux7eedux80a1ux606fux7387ux4f1aux600eux6837}

\textbf{问题。}纯作条件模型，取看涨3.96711451102和看跌0.4，使含股息平价\(C-P=S_0e^{-q\tau}-Ke^{-r\tau}\)成立的\(q\)是多少？

\textbf{解。}先解指数，
\[\boxed{q=-\frac1\tau\log\frac{C-P+Ke^{-r\tau}}{S_0}\approx0.0405803506315.}\]
这只说明加入这样一个假设可以配平价差，不证明真实公司支付这个股息率，也不证明同一波动率的两个含股息BS价格同时匹配。不能把反解出的参数冒作外部事实。

\bigskip

\section{习题6.25：用平价检查另一组历史输入}\label{ux4e60ux98986.25ux7528ux5e73ux4ef7ux68c0ux67e5ux53e6ux4e00ux7ec4ux5386ux53f2ux8f93ux5165}

\textbf{书内依据：第6.6节连续模型，第6.3节定理6.6。原题印刷250页。}给定\(S_0=81.63,K=70,\tau=0.3123,r=0.01,\sigma=0.26\)；看涨报价12.6---12.7，看跌1.43。

\subsection{（a）计算看涨}\label{aux8ba1ux7b97ux770bux6da8}

\[\sigma\sqrt\tau\approx0.14529790088,\quad
Ke^{-r\tau}\approx69.781731004437,\]
\[d_1\approx1.15198043744,\qquad d_2\approx1.00668253656.\]
代入得\(\boxed{C\approx12.630681006764}\)，处于指定区间内。

\subsection{（b）计算看跌并检查是否一致}\label{bux8ba1ux7b97ux770bux8dccux5e76ux68c0ux67e5ux662fux5426ux4e00ux81f4}

\[\boxed{P=C+69.781731004437-81.63\approx0.782412011201.}\]
这不同于1.43。只依据看涨区间，平价允许的看跌区间也仅为
\[\boxed{[0.751731004437,\ 0.851731004437]}. \]
因此给定看跌报价不在同一无分红理想模型允许的范围。不能因为一个看涨价格落在区间内，就称整组输入均合理。

\subsection{6.25-A：已知看跌确实产生正支付}\label{aux5df2ux77e5ux770bux8dccux786eux5b9eux4ea7ux751fux6b63ux652fux4ed8}

\textbf{问题。}在风险中性分布下，条件\(S_\tau<70\)后，平均到期看跌支付是多少？

\textbf{解。}正支付概率为\(\Phi(-d_2)\)，未折现支付均值为\(e^{r\tau}P\)。由于零支付不贡献期望，条件均值为
\[\boxed{\frac{e^{r\tau}P}{\Phi(-d_2)}\approx4.99771344623.}\]
这个条件均值高于无条件均值，并非看跌的当前公允价格，也不是实际世界的预测均值。

\subsection{6.25-B：可否只改波动率同时匹配}\label{bux53efux5426ux53eaux6539ux6ce2ux52a8ux7387ux540cux65f6ux5339ux914d}

\textbf{问题。}固定看涨12.65、看跌1.43和原股票、利率、期限，是否存在某个波动率使两价都匹配？

\textbf{解。}两报价差为\(12.65-1.43=11.22\)；无分红平价却要求\(S_0-PV(K)=11.848268995563\)。两数不等，而该恒等式不含波动率，所以\(\boxed{\text{不存在}}\)。先做代数一致性检查，比盲目对波动率数值求根更基本。

\subsection{6.25-C：看跌最多赔20}\label{cux770bux8dccux6700ux591aux8d5420}

\textbf{问题。}求\(\min\{(70-S_\tau)^+,20\}\)价格。

\textbf{解。}逐终值有\(\min\{(70-y)^+,20\}=(70-y)^+-(50-y)^+\)。用相同参数、执行价50再算看跌\(P_{50}\approx0.000816787010109\)，两份相减得
\[\boxed{P_{\rm cap}=P_{70}-P_{50}\approx0.781595224191065.}\]
这是一条精确支付恒等式加数值积分，不是把原价直接取\(\min(P,20)\)。

\clearpage

\section{本批验收清单与接续状态}\label{ux672cux6279ux9a8cux6536ux6e05ux5355ux4e0eux63a5ux7eedux72b6ux6001}

\begin{longtable}[]{@{}lrlr@{}}
\toprule\noalign{}
范围 & 原题数量 & 解答内容 & 变式数量 \\
\midrule\noalign{}
\endhead
\bottomrule\noalign{}
\endlastfoot
5.1---5.10 & 10 & 全部小问、鞅可积性、停止与极限条件 & 30 \\
6.1---6.5 & 5 & 一期复制、完整三状态线性方程、端点套利 & 15 \\
6.6---6.15 & 10 & 全部有限树价值和复制；完整终端支付 & 30 \\
6.16---6.20 & 5 & 美式全节点比较、最优停止和两种6.20解释 & 15 \\
6.21---6.25 & 5 & 方程验证、积分定价、平价和历史输入检查 & 15 \\
合计 & 35 & 原题与变式分开统计 & 105 \\
\end{longtable}

表中``完成''包括对不一致条件作明确的分情况回答，不包括为错误陈述强行制造肯定证明。5.6严格越界不能取其声称等号；6.3区间端点在通常套利定义下须排除；6.7按支付式与文字给不同合约；6.20两种支付都解，而不决定作者原意。核对美式6.17时，第二期HH的立即支付为18、继续价值21，正确决策是继续；不能把21误当成立即支付。

到本批为止，Durrett这一指定版本第5章10题、第6章25题均按顺序处理；第6章最后题号为6.25，随后进入附录。前四章仍保留此前明确列出的原版页面与题面问题，不能把本批完成升级成整本所有来源均已最终核清，更不能算作其余四本教材已完成。历史重复稿件不在这里重新累计。

\subsection{数值检查与数学证明分别是什么}\label{ux6570ux503cux68c0ux67e5ux4e0eux6570ux5b66ux8bc1ux660eux5206ux522bux662fux4ec0ux4e48}

本批程序以有理数重建所有有限树，逐节点核对银行与股票组合在上涨、下跌两个分支的支付；美式树还枚举有限停止策略复算最优值。正态公式另以原始支付积分复算，停止公式检查有限区间、有限路径与符号恒等式。题号、全部A/B/C标题及关键条件标记也单独检查。

检查不会证明没有拿到的2016原始页图，也不是对所有无限模型证明的形式认证或独立同行审稿。最终实际执行数量、PDF页数和版面检查记录见随附JSON与文本日志；不沿用未保存的中间运行数字。

\section{来源与书内定位}\label{ux6765ux6e90ux4e0eux4e66ux5185ux5b9aux4f4d}

{[}S1{]} Richard Durrett, \emph{Essentials of Stochastic Processes},
third edition, Springer, 2016,
ISBN9783319456133，DOI10.1007/978-3-319-45614-0。第5.5节原题印刷220---222页；第6.7节原题247---250页。可检索原书转录（不是原始页图）：

\url{https://dokumen.pub/essentials-of-stochastic-processes-3ed-331945613x-978-3-319-45613-3-978-3-319-45614-0-237-241-244-2.html}

{[}S2{]} 作者公开稿 \emph{Essentials of Stochastic Processes},
Version3.9,
May2021。第5章习题印刷182---183页，第6章相关习题印刷204---206页；用于核对对应模型、公式与明确变化，不把同号默认同题。

\url{https://sites.math.duke.edu/~rtd/EOSP/EOSP2021.pdf}

{[}S3{]}
作者2011年第二版公开稿。第5章相关题在印刷176---177页；第6章相关题印刷203---206页。2016的5.5可对应2011的5.13；2016的6.10---6.25依实际内容与此稿相应题核对。

\url{https://sites.math.duke.edu/~rtd/EOSP/EOSP2E.pdf}

作者入口：\url{https://sites.math.duke.edu/~rtd/EOSP/eosp.html}。本文所有解答和自编变式均为独立推导；随附资料包不包含出版社整书或字体文件。

\end{document}